\chapimg{fig/sh05.jpg}
\IfFileExists{extra/lead05.tex}{\chapterlead{\subhd{After this chapter you can} compute reference evapotranspiration by the Penman--Monteith equation; build a crop coefficient curve from the four growth stages; split evaporation into its two stages and compute the duration of stage 1; write the yield response to water stress; compute the net and gross irrigation requirement of a season; and read the crop library of the application.}
}{}
\chapter{Evapotranspiration, Crop Stress and the Crop Library}\label{chap:05}
\markboth{Evapotranspiration, Crop Stress and the Crop Library}{Evapotranspiration, Crop Stress and the Crop Library}
\IfFileExists{extra/bg05.tex}{\section{Introduction: the water that a crop uses}
A growing crop returns most of the water that reaches its root zone to the atmosphere, by \kw{transpiration} through its leaves and by \kw{evaporation} from the soil and from wet surfaces; the two together are \kw{evapotranspiration} (ET). ET is the largest outflow of the water balance of Chapter~2, it fixes the quantity of water that an irrigation system must supply, and it is the flux that a mulch or a trigger can change. It is driven by energy (the sun evaporates the water) and limited by supply (a dry soil cannot yield it). The standard method, developed by the FAO, separates these two aspects: the \emph{atmospheric demand} is expressed by a reference evapotranspiration $ET_0$ computed from weather alone; the \emph{crop} is expressed by a coefficient $K_c$ that scales $ET_0$ to its own demand $ET_c=K_cET_0$; and the \emph{stress} of a dry soil is a further factor $K_s\le1$ so that the actual rate is $ET_a=K_sK_cET_0$.

\section{Reference evapotranspiration: the Penman--Monteith equation}
The reference crop is a hypothetical grass of 0.12 m, with a fixed surface resistance of 70 s\,m$^{-1}$ and an albedo of 0.23, well watered and covering the ground. The FAO Penman--Monteith equation for a daily time step is

\begin{equation}ET_0=\frac{0.408\,\Delta\,(R_n-G)+\gamma\dfrac{900}{T+273}\,u_2\,(e_s-e_a)}{\Delta+\gamma\,(1+0.34\,u_2)},\label{eq:pm}\end{equation}


where $R_n$ is the net radiation (MJ\,m$^{-2}$\,d$^{-1}$), $G$ the soil heat flux (about zero daily), $T$ the mean temperature ($^\circ$C), $u_2$ the wind speed at 2 m (m\,s$^{-1}$), $e_s-e_a$ the vapour-pressure deficit (kPa), $\Delta$ the slope of the saturation vapour-pressure curve (kPa\,$^\circ$C$^{-1}$) and $\gamma$ the psychrometric constant (kPa\,$^\circ$C$^{-1}$). The first term of the numerator is the \emph{radiation} term, the second the \emph{aerodynamic} term; in humid climates the first dominates, in arid windy ones the second. Simpler equations (Hargreaves, Priestley--Taylor, pan evaporation with a pan coefficient) are used when data are missing; they are calibrated against Penman--Monteith.

\exbox{Penman--Monteith for a summer day}{Compute $ET_0$ for a summer day at 100 m altitude: $T_{max}=31$, $T_{min}=17$ $^\circ$C ($T=24$), $R_n=14.5$ MJ\,m$^{-2}$\,d$^{-1}$, $u_2=2.0$ m\,s$^{-1}$, mean relative humidity 55 \%.}\sol \givenline{$T_{max}=31$ ; $T_{min}=17$ $^\circ$C ; $T=24$ ; $R_n=14.5$ MJ\,m$^{-2}$\,d$^{-1}$ ; $u_2=2.0$ m\,s$^{-1}$}\ans Saturation vapour pressures: $e^\circ(31)=4.493$, $e^\circ(17)=1.938$, so $e_s=3.215$ kPa and $e_a=0.55e_s=1.768$ kPa; deficit $1.447$ kPa. $\Delta=4098\times0.6108\,e^{17.27\cdot24/(24+237.3)}/(24+237.3)^2=0.1791$. Pressure $P=101.3[(293-0.0065\times100)/293]^{5.26}=100.1$ kPa, $\gamma=0.000665P=0.0666$. Numerator: $0.408\times0.1791\times14.5=1.060$; $\gamma(900/297)\times2.0\times1.447=0.584$; sum $1.643$. Denominator: $0.1791+0.0666(1+0.68)=0.2910$. $ET_0=5.65$ mm\,d$^{-1}$ \ans The radiation term gives 3.6 mm and the aerodynamic term 2.0 mm of the total. \finans

\section{The crop coefficient and the four growth stages}
The crop coefficient $K_c$ changes through the season with the cover of the canopy. The season is divided into four stages: \emph{initial} (germination and early growth, length $L_1$, $K_c=K_{c,ini}$ nearly constant), \emph{development} ($L_2$, $K_c$ rises linearly to $K_{c,mid}$), \emph{mid-season} ($L_3$, $K_c=K_{c,mid}$) and \emph{late season} ($L_4$, $K_c$ falls linearly to $K_{c,end}$). The crop library of the application holds $K_{c,ini}$, $K_{c,mid}$, $K_{c,end}$, the four lengths, the maximum height $h$, the rooting depth $R_{max}$, the depletion fraction $p$ and the yield-response factor $K_y$ for 48 crops. the table below shows six of them.

\begin{table}[!htbp]\centering\small\caption{Six crops of the library: coefficients, stage lengths, height, rooting depth, depletion fraction and yield-response factor.}
\begin{tabular}{lrrrlrrrr}\toprule
\hd{Crop} & \hd{$K_{c,ini}$} & \hd{$K_{c,mid}$} & \hd{$K_{c,end}$} & \hd{Stage lengths (d)} & \hd{$h$ (m)} & \hd{$R_{max}$ (cm)} & \hd{$p$} & \hd{$K_y$}\\\midrule
Maize & 0.3 & 1.2 & 0.6 & 15+30+40+20 = 105 & 1.5 & 135 & 0.5 & 1.28\\
Wheat & 0.4 & 1.15 & 0.41 & 20+30+45+25 = 120 & 0.9 & 150 & 0.5 & 1.05\\
Tomato & 0.6 & 1.1 & 1.0 & 30+40+40+25 = 135 & 0.6 & 120 & 0.4 & 1.2\\
Potato & 0.5 & 1.1 & 0.4 & 25+30+45+30 = 130 & 0.6 & 60 & 0.4 & 1.1\\
Onion & 0.7 & 1.05 & 0.7 & 15+25+70+40 = 150 & 0.5 & 60 & 0.3 & 1.18\\
Soybean & 0.4 & 1.1 & 0.5 & 20+30+45+25 = 120 & 0.8 & 95 & 0.5 & 1.07\\\bottomrule\end{tabular}\end{table}

\exbox{Maize coefficient through the season}{For maize (the table below), compute $K_c$ on days 10, 30, 60, 90 and 110 after sowing and $ET_c$ if $ET_0=5.5$ mm\,d$^{-1}$ on each of those days.}\sol \givenline{$ET_0=5.5$ mm\,d$^{-1}$}\ans day 10: $K_c=0.300$, $ET_c=1.65$ mm; day 30: $K_c=0.750$, $ET_c=4.12$ mm; day 60: $K_c=1.200$, $ET_c=6.60$ mm; day 90: $K_c=1.050$, $ET_c=5.78$ mm; day 110: $K_c=0.600$, $ET_c=3.30$ mm \ans Day 10 is in the initial stage ($K_c=0.30$); day 30 in development (interpolated between 0.30 at day 15 and 1.20 at day 45: $0.30+0.90\times15/30=0.75$); 60 and 90 are in mid-season (1.20) and 110 is in late season: $1.20-0.60\times(110-85)/20$ is clipped at the end value 0.60 after day 105. \finans

\section{Evaporation from the soil: stages 1 and 2}
Soil evaporation follows two stages after a wetting (Ritchie, 1972). In \kw{stage 1} the surface is wet and the evaporation is limited by energy: it proceeds at the potential rate $E_p$ until the cumulative evaporation reaches $U$, the stage-1 capacity (about 5--12 mm, small in sand and large in clay). In \kw{stage 2} the surface layer has dried, evaporation is limited by the conductivity of the soil and falls as the square root of the time since stage 1 ended: $E(t)=\beta\,[\sqrt{t}-\sqrt{t-1}]$ per day, with $\beta$ a soil constant. The duration of stage 1 is $t_1=U/E_p$. A mulch lowers $E_p$ and hence lengthens stage 1 but also reduces the loss of stage 2; Chapter~10 treats it in detail. The two coefficients $K_e$ (evaporation) and $K_{cb}$ (basal, transpiration) of the dual-$K_c$ method split $K_c$ accordingly.

\exbox{Duration of stage 1}{The stage-1 capacity of a loam is $U=9$ mm and the potential evaporation is $E_p=4.5$ mm\,d$^{-1}$ in a hot week and 1.8 mm\,d$^{-1}$ in a cool one. Find the length of stage 1 in each, and the cumulative evaporation after 10 days in the hot week if $\beta=3.5$ mm\,d$^{-1/2}$.}\sol \givenline{$U=9$ mm ; $E_p=4.5$ mm\,d$^{-1}$ ; $\beta=3.5$ mm\,d}\ans $t_1=U/E_p$: hot week $9/4.5=2.0$ d; cool week $9/1.8=5.0$ d \ans After 10 days in the hot week: stage 1 gave 9 mm in 2 d; stage 2 lasts 8 d and gives $\beta\sqrt{8}=3.5\times2.83=9.9$ mm; total $=18.9$ mm \finans

\section{Water stress and yield response}
When the soil is drier than the allowed depletion the crop closes its stomata, transpires less than the demand and yields less. The FAO relation (Doorenbos and Kassam) relates the relative yield loss to the relative evapotranspiration deficit:

\begin{equation}1-\frac{Y_a}{Y_m}=K_y\Bigl(1-\frac{ET_a}{ET_c}\Bigr),\label{eq:ky}\end{equation}


with $K_y$ the yield-response factor: above 1 for crops that are sensitive to water (maize, 1.28; tomato, 1.20), below 1 for tolerant ones (sorghum, 0.90; groundnut, 0.70). The stress coefficient, within the daily model, is $K_s=(TAW-D_r)/[(1-p)TAW]$ once $D_r>p\,TAW=RAW$, and 1 before.

\exbox{Yield loss from a deficit}{A maize crop with $ET_c=520$ mm in the season gets $ET_a=430$ mm because irrigation was short. Find the relative yield and the yield if $Y_m=11.5$ Mg\,ha$^{-1}$. What is the yield of a tolerant crop with $K_y=0.7$ for the same deficit?}\sol \givenline{$ET_c=520$ mm ; $ET_a=430$ mm because irrigation was s ; $Y_m=11.5$ Mg\,ha$^{-1}$ ; $K_y=0.7$}\ans Relative deficit $1-430/520=0.173$. Maize: $1-Y_a/Y_m=1.28\times0.173=0.222$, so $Y_a/Y_m=0.778$ and $Y_a=8.95$ Mg\,ha$^{-1}$ \ans Tolerant crop: $1-0.7\times0.173=0.879$, so a loss of 12.1 per cent against 22.2 per cent for maize. \finans

}{}
\banner{Evapotranspiration and the two stages of soil evaporation}
\section{The energy that drives evaporation}
Evaporation of one kilogram of water needs the latent heat $\lambda\approx2.45$ MJ\,kg$^{-1}$. A net radiation of $R_n=13$ MJ\,m$^{-2}$\,d$^{-1}$ would evaporate $13/2.45=5.3$ mm of water per day if all of it were used for evaporation. This is why potential evapotranspiration is called \emph{energy based}. The equivalent depth of evaporation is $R_n/\lambda=0.408\,R_n$ in mm\,d$^{-1}$ with $R_n$ in MJ\,m$^{-2}$\,d$^{-1}$.

\section{The Penman--Monteith equation}
For a well-watered grass reference, the FAO-56 form of the Penman--Monteith equation gives the reference evapotranspiration
\begin{equation}\label{eq:pm}\boxed{ET_0=\frac{0.408\,\Delta(R_n-G)+\gamma\,\dfrac{900}{T+273}\,u_2\,(e_s-e_a)}{\Delta+\gamma(1+0.34\,u_2)},}\end{equation}
\par\noindent where $ET_0$ is the reference evapotranspiration (mm\,d$^{-1}$), $\Delta$ the slope of the saturation vapour-pressure curve (kPa\,K$^{-1}$), $R_n$ the net radiation and $G$ the soil heat flux (MJ\,m$^{-2}$\,d$^{-1}$), $\gamma$ the psychrometric constant (kPa\,K$^{-1}$), $T$ the mean air temperature (\textdegree C), $u_2$ the wind speed at 2 m (m\,s$^{-1}$), $e_s$ the saturation and $e_a$ the actual vapour pressure (kPa); the constants 0.408, 900 and 0.34 follow from the units and from the reference grass surface.

in mm\,d$^{-1}$, where $\Delta$ is the slope of the saturation vapour pressure curve (kPa\,$^\circ$C$^{-1}$), $\gamma$ the psychrometric constant, $G$ the soil heat flux, $u_2$ the wind speed at 2 m, and $e_s-e_a$ the vapour pressure deficit. The first term is the radiation term and the second the aerodynamic term.
\exbox{A mild day}{Take $T=20\,^\circ$C, $R_n=13$ MJ\,m$^{-2}$\,d$^{-1}$, $G=0$, $u_2=2$ m\,s$^{-1}$, relative humidity 60 per cent and pressure 101.3 kPa. Then $e_s=0.6108\exp[17.27(20)/(20+237.3)]=2.338$ kPa, $e_a=0.60(2.338)=1.403$ kPa, $\Delta=4098\,e_s/(T+237.3)^2=0.1447$, and $\gamma=0.000665(101.3)=0.0674$. The radiation term is $0.408(0.1447)(13)=0.768$; the aerodynamic term is $0.0674\,(900/293)(2)(0.935)=0.387$; the denominator is $0.1447+0.0674(1+0.68)=0.2579$. So $ET_0=(0.768+0.387)/0.2579=4.48$ mm\,d$^{-1}$. About two thirds of it is driven by radiation.}



\section{Evaporation from bare soil: two stages}
Soil evaporation does not follow the atmospheric demand for long. In \emph{stage 1} the surface is wet and the soil evaporates at the potential rate, limited only by energy, until the cumulative evaporation reaches a limit $U$. In \emph{stage 2} the surface layer has dried, water reaches the surface as vapour or liquid through a dry layer, and the rate falls. A simple description (Ritchie, 1972) gives the cumulative evaporation in stage 2 as proportional to the square root of time since the start of the stage,
\begin{equation}\label{eq:stage2}\boxed{E(t)=U+\beta\sqrt{t-t_1},}\end{equation}
where $t_1$ is the time at which stage 1 ends and $\beta$ (mm\,d$^{-1/2}$) is a soil constant.
\exbox{Evaporation after a wetting}{With $U=9$ mm and $\beta=3.5$ mm\,d$^{-1/2}$ (typical textbook values for a loam) and a potential rate of 4.5 mm\,d$^{-1}$, stage 1 lasts $9/4.5=2$ days. After a further 1, 4 and 9 days the stage-2 cumulative evaporation is $3.5$, $7.0$ and $10.5$ mm: the rate has fallen from 4.5 to $\beta/(2\sqrt{t})=1.75$ mm\,d$^{-1}$ after one day and to 0.58 mm\,d$^{-1}$ after nine. Rain that wets the soil restarts stage 1, which is why small frequent wettings, from rain or from irrigation, lose a large share to evaporation: this is the 60 to 72 per cent loss that the ensemble finds for a 1 mm dose.}


\fig{sc_e1k}{0.75}{The SWAFL evaporation layer: from the pore to the atmosphere, with the two stages.}

\section{What a mulch changes}
A mulch lowers the energy that reaches the soil surface and the vapour exchange with the atmosphere, so it lowers $E_s$, mainly in stage 1, and it intercepts some rain. It does not change transpiration directly. Its effect on the balance is a change of $E_s$ of, say, $-\Delta E$; by Equation~(\ref{eq:bal}) this must be balanced by an increase of $T$, of $D$, of storage or by a decrease of irrigation. Chapter~7 gives the memory fate that tells how it is shared.

\begin{daisybox}{potential evapotranspiration in Daisy}
Daisy computes an energy-based potential evapotranspiration from the weather of its input file and divides it between the canopy and the soil according to the cover. The soil evaporation of Daisy then follows the water content of the top layers through the Richards equation, so the two-stage behaviour of this chapter appears in the output without being imposed. Daisy's logs give $T$, $E_s$ and the canopy evaporation separately, and the audit of Chapter~8 asks whether they add up to $ET$.
\end{daisybox}

\exbox{}{Compute $ET_0$ for $T=30\,^\circ$C, $R_n=18$, $G=0$, $u_2=3$, relative humidity 35 per cent, $P=101.3$ kPa.}\sol  $e_s=4.243$, $e_a=1.485$, $\Delta=0.2434$, $\gamma=0.0674$; $ET_0=9.07$ mm\,d$^{-1}$.

\exbox{}{Derive the equivalent evaporation depth $R_n/\lambda$ in mm for $R_n=1$ MJ\,m$^{-2}$\,d$^{-1}$ using $\lambda=2.45$ MJ\,kg$^{-1}$ and $\rho_w=1000$ kg\,m$^{-3}$.}\sol  $1/2.45=0.408$ mm per MJ\,m$^{-2}$\,d$^{-1}$.

\exbox{}{With $U=6$ mm, $\beta=2.5$ mm\,d$^{-1/2}$ and a potential rate of 5 mm\,d$^{-1}$, find the cumulative evaporation 6 days after a wetting.}\sol  Stage 1 lasts $6/5=1.2$ d; $t-t_1=4.8$ d; $E=6+2.5\sqrt{4.8}=11.5$ mm.

\exbox{}{Explain why two 5 mm irrigations one week apart can lose more to evaporation than one 10 mm irrigation.}\sol  Each wetting restarts stage 1, where evaporation proceeds at the potential rate; a single large wetting spends most of its time in stage 2, with the dry surface layer slowing evaporation.
\banner{Crop water stress and yield}
\section{Available water and depletion}
The total available water of a root zone is $TAW=1000(\theta_{fc}-\theta_{wp})Z_r$ in mm, with $Z_r$ the rooting depth in m. In the ensemble the total available water to 1 m runs from 60 mm (sand) to 220 mm (silt loam) (Table~\ref{tab:soils}). The depletion $D_r$ is the water missing from field capacity. A crop does not suffer until the depletion exceeds a fraction $p$ of $TAW$, the readily available water $RAW=p\,TAW$.

\section{The stress coefficient}
Beyond $RAW$ the transpiration falls linearly with the depletion:
\begin{equation}\label{eq:Ks}\boxed{K_s=\begin{cases}1,&D_r\le RAW\\[2pt] \dfrac{TAW-D_r}{(1-p)\,TAW},&D_r>RAW\end{cases}\qquad ET_a=K_s\,K_c\,ET_0 .}\end{equation}
\par\noindent where $K_s$ is the water-stress coefficient (0 to 1), $D_r$ the root-zone depletion (mm), ${\rm RAW}=p\,{\rm TAW}$ the readily available water (mm), ${\rm TAW}$ the total available water (mm), $p$ the fraction of TAW that the crop can use without stress, $ET_a$ the actual and $ET_0$ the reference evapotranspiration (mm\,d$^{-1}$) and $K_c$ the crop coefficient.

\exbox{A loam under drying}{The loam S06 has $TAW=150$ mm; take $p=0.5$, so $RAW=75$ mm. With $K_c=1.15$ and $ET_0=4.48$ mm\,d$^{-1}$ (Chapter~13) the unstressed demand is $ET_c=5.15$ mm\,d$^{-1}$. At $D_r=60$ and 75 mm, $K_s=1$ and $ET_a=5.15$. At $D_r=100$ mm, $K_s=(150-100)/(0.5\cdot150)=0.667$ and $ET_a=3.43$ mm\,d$^{-1}$. At $D_r=120$, $K_s=0.40$ and $ET_a=2.06$. The depletion at which the field must be irrigated to avoid stress is $RAW=75$ mm; irrigating at the trigger gives the schedule of the application.}



\section{From deficit to yield}
A widely used linear relation (FAO-33) links the relative yield loss to the relative evapotranspiration deficit:
\begin{equation}\label{eq:Ky}\boxed{1-\frac{Y_a}{Y_m}=K_y\Bigl(1-\frac{ET_a}{ET_m}\Bigr).}\end{equation}
\par\noindent where $Y_a$ and $Y_m$ are the actual and maximum yields (Mg\,ha$^{-1}$), $ET_a$ and $ET_m$ the actual and maximum seasonal evapotranspiration (mm) and $K_y$ the yield response factor of the crop (dimensionless).

\exbox{A hypothetical cereal}{With a yield response factor $K_y=1.15$ (illustrative) and a season in which $ET_a/ET_m=0.80$, the relative yield is $Y_a/Y_m=1-1.15(0.20)=0.77$. If $Y_m=6.6$ Mg\,ha$^{-1}$, as in the best triggers of the ensemble (the reach figure of Chapter~5, panel d), $Y_a=5.1$ Mg\,ha$^{-1}$. The ensemble's rainfed sand yields about 1.1 Mg\,ha$^{-1}$, because a trigger beyond the reach means no irrigation at all.}



\section{Link with the trigger}
The suction at the sensor at which the depletion reaches $RAW$ is the \emph{ideal} trigger. If it lies beyond the reach $\hc$ of the soil, the crop cannot reach it and the field is stressed before the sensor says so. Comparing the two is the use of Chapter~5 in design: choose the trigger below the smaller of the two, and below $\hc/2$ in practice.
\exbox{}{A soil has $TAW=95$ mm (S05) and $p=0.55$. Find $RAW$ and $K_s$ at $D_r=70$ mm.}\sol  $RAW=0.55(95)=52.3$ mm; $K_s=(95-70)/(0.45\cdot95)=0.585$.

\exbox{}{With $K_y=0.9$ and a season mean $ET_a/ET_m=0.7$, find the relative yield.}\sol  $1-0.9(0.3)=0.73$.

\exbox{}{For sand (S03, $TAW=60$ mm, $p=0.5$) find how many days at $ET_c=5$ mm\,d$^{-1}$ pass between a full profile and the onset of stress. Compare with the loam.}\sol  Sand: $RAW=30$ mm, $30/5=6$ d. Loam: $RAW=75$ mm, $15$ d.

\exbox{}{Explain why a sensor trigger beyond the reach leads to a larger yield loss in a sand than the stress coefficient alone would predict.}\sol  A sensor trigger beyond the reach is never read, so the crop is not irrigated even as the stress coefficient falls; the field is rainfed and the sand's small $TAW$ means the stress comes within days.
\banner{Crop library and crop water requirement}\label{ch:crops}
\textit{Illustrative climate used throughout.} The chapter uses a synthetic year so that every number can be reproduced: reference evapotranspiration $ET_0=4.0+2.0\sin[2\pi(d-100)/365]$ mm d$^{-1}$ with a small seeded noise (mean 4 mm d$^{-1}$, 1,460 mm per year), and rain on 25\% of the days with exponential depths of mean 8.2 mm (about 600 mm per year; seed 7). The soil is the loam with $\mathrm{AWC}=1.7$ mm cm$^{-1}$ and the trigger suction $-300$ cm, which corresponds to the depletion fraction $f_T=0.468$ of total available water (Eq.~5.3). Planting is on day 100.

\section{The crop-coefficient curve}
\begin{equation}K_c(t)=\begin{cases}K_{c,\rm ini}&t<L_1\\K_{c,\rm ini}+(K_{c,\rm mid}-K_{c,\rm ini})\dfrac{t-L_1}{L_2}&L_1\le t<L_1+L_2\\K_{c,\rm mid}&L_1+L_2\le t<L_1+L_2+L_3\\K_{c,\rm mid}+(K_{c,\rm end}-K_{c,\rm mid})\min\!\Big(1,\dfrac{t-L_1-L_2-L_3}{L_4}\Big)&\text{after}\end{cases}\end{equation}
with $t$ days after planting and stage lengths $L_1\ldots L_4$ (initial, development, mid-season, late). Outside the season the bare soil has $K_c=0.25$ (the FAO-56 single-coefficient curve).

\exbox{Reading the curve}{Maize: $K_{c,\rm ini}=0.30$, $K_{c,\rm mid}=1.20$, $K_{c,\rm end}=0.60$; stages $15/30/40/20$ d. Find $K_c$ at $t=10,\,30,\,60,\,95$ d.}\sol  $t=10<15$: $K_c=0.30$. $t=30$ lies in the development stage ($15\le30<45$): $K_c=0.30+0.90\times(30-15)/30=0.30+0.45=0.75$. $t=60$ lies in the mid-season ($45\le60<85$): $K_c=1.20$. $t=95$ is in the late season ($85\le95<105$): $K_c=1.20+(0.60-1.20)\times(95-85)/20=1.20-0.30=0.90$.


\exbox{Daily crop need}{The daily need is $ET_c=K_cET_0$. On a day with $ET_0=5.2$ mm d$^{-1}$ find $ET_c$ of maize at $t=30$ and $t=60$, and the total need of the development stage if $ET_0$ stays at 5.2.}\sol  $t=30$: $0.75\times5.2=3.90$ mm. $t=60$: $1.20\times5.2=6.24$ mm. The development stage covers $t=15\ldots44$, $K_c$ rising linearly from 0.30 to 1.20; its mean is 0.75 (the average of the end points), so the need is $30\times0.75\times5.2=117.0$ mm.


\section{Roots, total and readily available water, and stress}
The roots of an annual crop grow with the canopy,
\begin{equation}\boxed{z_r=\mathrm{clip}\Big(15+(R_{\max}-15)\min\Big(1,\frac{1.15K_c}{\max(0.3,K_{c,\rm mid})}\Big);\,15,R_{\max}\Big)\ \text{cm},\quad R_{\max}\le100,}\end{equation}
\par\noindent where $z_r$ is the rooting depth (cm), $R_{\max}$ the maximum depth allowed by the soil and the crop (cm), $K_c$ the crop coefficient of the day, $K_{c,{\rm mid}}$ its mid-season value, and ${\rm clip}(x;a,b)$ limits $x$ to the interval $[a,b]$.

and the total and readily available water follow FAO-56:
\begin{equation}\mathrm{TAW}=\mathrm{AWC}\,z_r,\quad p_d=\mathrm{clip}\{p+0.04(5-K_cET_0);0.1,0.8\},\quad \mathrm{RAW}=p_d\mathrm{TAW},\quad K_s=\begin{cases}1&D_r\le\mathrm{RAW}\\\dfrac{\mathrm{TAW}-D_r}{(1-p_d)\mathrm{TAW}}&D_r>\mathrm{RAW}\end{cases}\end{equation}
\par\noindent where ${\rm AWC}$ is the available water capacity of the soil (mm\,cm$^{-1}$), $p_d$ the daily depletion fraction (adjusted for the demand $K_cET_0$ in mm\,d$^{-1}$ and kept between 0.1 and 0.8), $p$ the tabulated fraction, RAW and TAW as before, $D_r$ the depletion and $K_s$ the stress coefficient.

and $ET_a=K_sET_c$. The depletion $D_r$ changes by $D_{r,d}=D_{r,d-1}-(P_d-RO_d)+ET_{a,d}$ with $RO_d=0.3\max(0,P_d-35)$, and a negative $D_r$ drains below the roots. Irrigation refills the root zone to field capacity on the day the depletion reaches the trigger, $D_r\ge f_T\mathrm{TAW}$: the net depth is $I_n=D_r$ and the gross depth $I=I_n/E_{a,m}$ with $E_{a,m}=0.60$ (surface), 0.75 (overhead), 0.90 (drip).

\exbox{TAW, RAW and the stress coefficient}{Maize ($R_{\max}=135$ cm, capped at 100 cm by the model profile; $p=0.50$) on the loam (AWC $=1.7$ mm cm$^{-1}$). At $t=30$ ($K_c=0.75$) with $ET_0=5.2$ mm d$^{-1}$, find $z_r$, TAW, $p_d$, RAW, and $K_s$ when $D_r=40$ mm and when $D_r=55$ mm.}\sol  $1.15K_c/\max(0.3,1.20)=1.15\times0.75/1.20=0.719$; $z_r=15+85\times0.719=76.1$ cm. TAW $=1.7\times76.1=129.4$ mm. $p_d=0.50+0.04(5-0.75\times5.2)=0.50+0.04\times1.1=0.544$; RAW $=0.544\times129.4=70.4$ mm. With $D_r=40\le$ RAW: $K_s=1$. With $D_r=55>$ RAW: $K_s=(129.4-55)/((1-0.544)\times129.4)=74.4/59.0=1.261$.


\begin{table}[!htbp]\centering\small\caption{Growth of the maize root zone on the loam ($ET_0=4.5$ mm d$^{-1}$ in $p_d$).}\label{tab:roots}
\begin{tabular}{lrrrrr}\toprule
\hd{$t$ (d)} & \hd{$K_c$} & \hd{$z_r$ (cm)} & \hd{TAW (mm)} & \hd{$p_d$} & \hd{RAW (mm)}\\\midrule
0 & 0.30 & 39.4 & 67.0 & 0.646 & 43.3\\
10 & 0.30 & 39.4 & 67.0 & 0.646 & 43.3\\
20 & 0.45 & 51.7 & 87.8 & 0.619 & 54.4\\
30 & 0.75 & 76.1 & 129.4 & 0.565 & 73.1\\
45 & 1.20 & 100.0 & 170.0 & 0.484 & 82.3\\
60 & 1.20 & 100.0 & 170.0 & 0.484 & 82.3\\
85 & 1.20 & 100.0 & 170.0 & 0.484 & 82.3\\
104 & 0.63 & 66.3 & 112.7 & 0.587 & 66.1\\\bottomrule\end{tabular}\end{table}

\section{A season of maize}
Running the balance day by day from planting gives the irrigations of the season. Table~\ref{tab:maizeirr} lists them for surface irrigation ($E_{a,m}=0.60$).

\begin{table}[!htbp]\centering\small\caption{Irrigations of maize on the loam, trigger $-300$ cm, surface irrigation.}\label{tab:maizeirr}
\begin{tabular}{lrrr}\toprule
\hd{day after planting} & \hd{day of year} & \hd{net $I_n$ (mm)} & \hd{gross $I$ (mm)}\\\midrule
42 & 142 & 81.5 & 135.8\\
55 & 155 & 81.6 & 136.1\\
68 & 168 & 81.1 & 135.2\\
84 & 184 & 84.0 & 140.0\\
104 & 204 & 60.3 & 100.4\\\bottomrule\end{tabular}\end{table}

\exbox{Four days to an irrigation}{Follow the four days up to the first irrigation of Table~\ref{tab:maizeirr}.}\sol  The depletion on the day before is $D_r=63.24$ mm. Each day: $D_r\leftarrow D_r-(P-RO)+ET_a$.
\begin{itemize}
\item day 39: $K_c=0.990$, $ET_0=5.16$, $ET_c=5.11$, $z_r=95.6$, TAW $=162.6$, RAW $=80.6$, $K_s=1.000$, $ET_a=5.11$, $P=0.0$, $D_r=68.35$ mm, $I_n=0.0$
\item day 40: $K_c=1.020$, $ET_0=4.90$, $ET_c=5.00$, $z_r=98.1$, TAW $=166.7$, RAW $=83.4$, $K_s=1.000$, $ET_a=5.00$, $P=0.0$, $D_r=73.35$ mm, $I_n=0.0$
\item day 41: $K_c=1.050$, $ET_0=5.66$, $ET_c=5.94$, $z_r=100.0$, TAW $=170.0$, RAW $=78.6$, $K_s=1.000$, $ET_a=5.94$, $P=3.3$, $D_r=75.96$ mm, $I_n=0.0$
\item day 42: $K_c=1.080$, $ET_0=5.11$, $ET_c=5.52$, $z_r=100.0$, TAW $=170.0$, RAW $=81.5$, $K_s=1.000$, $ET_a=5.52$, $P=0.0$, $D_r=0.00$ mm, $I_n=81.5$
\end{itemize}
The trigger depletion is $f_T\mathrm{TAW}=0.468\times170.0=79.5$ mm. On the last day $D_r$ reaches that value, so the layer irrigates the net depth $I_n=D_r=81.5$ mm and the root zone returns to field capacity ($D_r=0$). The gross depth for surface irrigation is $81.5/0.60=135.8$ mm.


\begin{table}[!htbp]\centering\small\caption{The same season under three methods.}\label{tab:method}
\begin{tabular}{lrrrrr}\toprule
\hd{Method} & \hd{$E_{a,m}$} & \hd{irrigations} & \hd{net (mm)} & \hd{gross (mm)} & \hd{lost (mm)}\\\midrule
surface & 0.60 & 5 & 388 & 647 & 259\\
overhead & 0.75 & 5 & 388 & 518 & 129\\
drip & 0.90 & 5 & 388 & 432 & 43\\\bottomrule\end{tabular}\end{table}

\exbox{What the method saves}{From Table~\ref{tab:method}, what does moving from surface to drip save in water per hectare?}\sol  Net need is the same in both (388 mm); gross falls from 647 mm to 432 mm, a saving of 216 mm $=2{,}158$ m$^3$ ha$^{-1}$ (1 mm over 1 ha is 10 m$^3$). The saving is the difference of the losses: $388(1/0.60-1/0.90)=216$ mm.


\section{Relative yield}
The relative yield of a season is the FAO-33 response
\begin{equation}\boxed{Y_{\rm rel}=\mathrm{clip}\Big(1-K_y\Big(1-\frac{\sum ET_a}{\sum ET_c}\Big);0,1\Big),}\end{equation}
with $K_y=1$ when the crop has no tabulated value.

\begin{table}[!htbp]\centering\small\caption{Maize on the loam at different triggers (surface irrigation, synthetic year).}\label{tab:trig}
\begin{tabular}{lrrrrrr}\toprule
\hd{trigger (cm)} & \hd{$f_T$} & \hd{irrigations} & \hd{gross (mm)} & \hd{$\sum ET_a$} & \hd{$\sum ET_a/\sum ET_c$} & \hd{$Y_{\rm rel}$}\\\midrule
50 & 0.000 & 105 & 734 & 514 & 1.000 & 1.000\\
100 & 0.000 & 105 & 734 & 514 & 1.000 & 1.000\\
200 & 0.321 & 7 & 629 & 514 & 1.000 & 1.000\\
300 & 0.468 & 5 & 647 & 513 & 0.998 & 0.997\\
500 & 0.614 & 3 & 524 & 482 & 0.939 & 0.922\\
800 & 0.718 & 3 & 556 & 473 & 0.920 & 0.898\\
1500 & 0.821 & 2 & 424 & 406 & 0.790 & 0.731\\
5000 & 0.941 & 1 & 267 & 364 & 0.708 & 0.626\\
50000 & 1.000 & 1 & 194 & 252 & 0.490 & 0.347\\\bottomrule\end{tabular}\end{table}

\exbox{Rainfed maize}{Compute $Y_{\rm rel}$ of maize with no irrigation, $K_y=1.28$.}\sol  With rain only, the season has $\sum ET_c=514$ mm and $\sum ET_a=252$ mm, so the ratio is $0.490$. $Y_{\rm rel}=1-1.28\times(1-0.490)=0.347$. A deficit of 51.0\% of the crop need becomes a yield loss of 65.3\%, because $K_y>1$ (maize is sensitive). A tolerant crop such as groundnut ($K_y=0.70$) would lose only 35.7\% for the same deficit.


\begin{keybox}Table~\ref{tab:trig} shows the same shape as the manuscript trigger experiment (Chapter~\ref{ch:e2}): yield is flat over a wide range of wet triggers, then falls as the trigger goes drier, while the gross irrigation falls overall (not strictly at every step, because a wetter trigger irrigates earlier and rain refills part of the root zone). At the wettest triggers ($f_T=0$) the rule irrigates every day, which is the model's limit of continuous irrigation. The numbers come from a synthetic climate, so they differ from the manuscript's; the shape is the point.\end{keybox}
\section{Every crop on one soil}
Table~\ref{tab:lib} runs the balance for twelve library crops of Table~21.2 of the application on the same loam, weather, trigger ($-300$ cm) and method (surface), planting on day 100.

\begin{table}[!htbp]\centering\small\caption{Twelve crops on the loam (synthetic year, surface method).}\label{tab:lib}
\begin{tabular}{lrrrrrrrrr}\toprule
\hd{Crop} & \hd{$K_{c,\rm ini}$} & \hd{$K_{c,\rm mid}$} & \hd{$K_{c,\rm end}$} & \hd{days} & \hd{$\sum ET_c$} & \hd{net} & \hd{irr.} & \hd{$K_y$} & \hd{$Y_{\rm rel}$}\\\midrule
Maize & 0.30 & 1.20 & 0.60 & 105 & 514 & 388 & 5 & 1.28 & 0.997\\
Wheat & 0.40 & 1.15 & 0.41 & 120 & 570 & 389 & 5 & 1.05 & 0.999\\
Sorghum & 0.30 & 1.05 & 0.45 & 125 & 527 & 323 & 4 & 0.90 & 1.000\\
Soybean & 0.40 & 1.10 & 0.50 & 120 & 556 & 379 & 5 & 1.07 & 0.999\\
Groundnut & 0.40 & 1.05 & 0.60 & 120 & 534 & 323 & 4 & 0.70 & 1.000\\
Cotton & 0.40 & 1.10 & 0.50 & 190 & 839 & 508 & 7 & 1.05 & 1.000\\
Tomato & 0.60 & 1.10 & 1.00 & 135 & 675 & 456 & 6 & 1.20 & 0.978\\
Potato & 0.50 & 1.10 & 0.40 & 130 & 602 & 405 & 9 & 1.10 & 0.986\\
Onion & 0.70 & 1.05 & 0.70 & 150 & 767 & 487 & 10 & 1.18 & 0.939\\
Cabbage & 0.70 & 1.05 & 1.00 & 165 & 796 & 516 & 13 & 0.95 & 0.981\\
Cassava & 0.30 & 1.00 & 0.60 & 210 & 876 & 561 & 11 & 1.00 & 1.000\\\bottomrule\end{tabular}\end{table}

\exbox{Comparing the crops}{Which crop needs the most water in Table~\ref{tab:lib}, which gets the most irrigations, and why can the two differ?}\sol  The largest seasonal need is Cassava ($\sum ET_c=876$ mm over 210 days) and the smallest is Maize ($514$ mm over 105 days). Need follows the length of the season and the height of $K_c$; the number of irrigations follows the size of the root-zone reservoir instead: Cabbage gets the most irrigations (13), because its small TAW (shallow roots, $R_{\max}$ = 50 cm) is used up quickly, even though its seasonal need is lower than the largest. The irrigation is also only the part of the need not met by rain and stored water, so it depends on when in the year the crop grows. The lowest relative yield is Onion ($Y_{\rm rel}=0.939$), where the shallow roots and the small $p$ make the root zone run dry between irrigations at this trigger.


\section{Carrying a library crop into the cereal-based model}
The annual model (Chapter~\ref{ch:e2}) is built from cereal runs. A library crop is transferred from it by the ratio of two runs of the balance above on the same soil, weather, trigger and method, one for the crop ($c$) and one for the cereal reference ($r$, small grains $K_c=0.30/1.15/0.25$). With $m$ the annual-model result of the cereal,
\begin{align}I&=\begin{cases}I_m\,\mathrm{clip}(I_c/I_r;0,4)&I_r>5\text{ mm}\\I_c&\text{otherwise}\end{cases}\\
T&=\min\Big(0.9ET,\;T_m\,\mathrm{clip}\Big(\frac{ET_{a,c}}{ET_{c,c}}\cdot\frac{ET_{c,c}}{ET_{c,r}}\Big(\frac{f_{c,c}}{f_{c,r}}\Big)^{0.3};0.2,3\Big)\Big),\qquad E^*=E_m\,\mathrm{clip}\Big(\frac{1-0.7f_{c,c}}{1-0.7f_{c,r}};0.6,1.6\Big)\\
D&=\max(0.2D_m,D_m+0.5(DP_c-DP_r)),\quad ET=\mathrm{clip}(P+I-D-\Delta S_m;50,5000),\quad N_L=N_{L,m}\frac{D+5}{D_m+5}.\end{align}
The remainder of the balance, $ET-T$, is shared between soil evaporation and the other evaporation terms in proportion to their targets ($E_{\rm other}^*=E_{{\rm other},m}ET/ET_m$).

\exbox{Transferring sorghum from the cereal result}{The cereal model run of the loam (trigger $-300$ cm, 600 mm rain; Chapter~\ref{ch:e2}) gives $P=606$, $I_m=90$, $D_m=200$, $E_m=158$, $T_m=230$ mm, so $ET_m=P+I_m-D_m=496$ mm and $E_{{\rm other},m}=ET_m-T_m-E_m=108$ mm; assume $\Delta S_m=0$ and $N_{L,m}=40$ kg ha$^{-1}$. Transfer sorghum with ground covers $f_{c,c}=0.80$ and $f_{c,r}=0.85$ (assumed mid-season values). The balances of this chapter give: sorghum $I_c=538$ mm, $ET_{a,c}/ET_{c,c}=1.000$, $ET_{c,c}=527$ mm, $DP_c=11$ mm; reference $I_r=609$ mm, $ET_{c,r}=543$ mm, $DP_r=0$ mm.}\sol  \emph{Irrigation:} $I=90\times\mathrm{clip}(538/609)=90\times0.884=80$ mm. \emph{Drainage:} $D=\max(0.2\times200,200+0.5(11-0))=206$ mm. \emph{Evapotranspiration:} $ET=P+I-D=606+80-206=480$ mm. \emph{Transpiration:} ratio $=1.000\times(527/543)\times(0.80/0.85)^{0.3}=0.952$, so $T=\min(0.9\times480,\,230\times0.952)=219$ mm. \emph{Evaporation:} $E^*=158\times\mathrm{clip}((1-0.56)/(1-0.595))=158\times1.086=172$ mm, $E^*_{\rm other}=108\times480/496=104$ mm; sharing $ET-T=261$ mm in that proportion gives $E=163$ and $E_{\rm other}=99$ mm. Leaching: $N_L=40\times(206+5)/(200+5)=41.1$ kg ha$^{-1}$. The balance closes: $T+E+E_{\rm other}=480=ET$.


\begin{keybox}These transfer rules are engineered, not measured: the clips (0.2--3, 0.6--1.6, 0--4) keep the result within the range in which the cereal simulation is informative, and the layer never gives a library crop the resolution of the cereal runs. The yield shown for a library crop is the relative yield of Equation 21.5, in per cent.\end{keybox}

\section{Summary of the crop layer}
\begin{itemize}
\item $K_c(t)$ in four stages; $ET_c=K_cET_0$; roots grow with the canopy.
\item TAW $=\mathrm{AWC}z_r$, RAW $=p_d$TAW with $p_d=p+0.04(5-K_cET_0)$; $K_s$ falls linearly below RAW.
\item The balance $D_r$ triggers irrigation at $f_T$TAW; gross $=$ net$/E_{a,m}$.
\item $Y_{\rm rel}=1-K_y(1-\sum ET_a/\sum ET_c)$.
\item A library crop enters the annual model by the ratio of its balance to the cereal reference.
\end{itemize}
\IfFileExists{extra/long05.tex}{\subsection{Complete calculation of evapotranspiration and stress}
\exbox{Reference ET, crop ET, stress and yield}{On a day with mean temperature 26$^\circ$C, relative humidity 45 per cent, net radiation 15.2 MJ\,m$^{-2}$\,d$^{-1}$, wind 2.1 m\,s$^{-1}$ at 2 m and altitude 100 m, find $ET_0$, then $ET_c$ and the actual ET of a crop at mid-season ($K_{cb}=1.10$, $K_e=0.10$) whose root zone holds $TAW=140$ mm and is depleted by 90 mm, and the resulting yield loss.}\sol
\ans Saturation vapour pressure $e_s=0.6108\exp[17.27T/(T+237.3)]=3.361$ kPa; actual $e_a=e_sRH/100=1.513$ kPa; deficit $D=1.849$ kPa
\ans Slope of the saturation curve $\Delta=4098e_s/(T+237.3)^2=0.1987$ kPa\,K$^{-1}$
\ans Pressure at 100 m: $P=101.3[(293-0.0065z)/293]^{5.26}=100.1$ kPa; psychrometric constant $\gamma=0.000665P=0.0666$ kPa\,K$^{-1}$
\ans Penman--Monteith: $ET_0=\dfrac{0.408\Delta(R_n-G)+\gamma\frac{900}{T+273}u_2D}{\Delta+\gamma(1+0.34u_2)}$ with $R_n=15.2$ MJ\,m$^{-2}$\,d$^{-1}$, $u_2=2.1$ m\,s$^{-1}$ gives $ET_0=6.43$ mm\,d$^{-1}$
\ans Dual crop coefficient at mid-season: $K_{cb}=1.10$, $K_e=0.10$, so $K_c=1.20$ and $ET_c=K_cET_0=7.71$ mm\,d$^{-1}$ without stress
\ans Stress: $TAW=140$ mm, $p=0.5$, depletion $D_r=90$ mm; $RAW=70$ mm; since $D_r>RAW$, $K_s=(TAW-D_r)/[(1-p)TAW]=0.714$
\ans Actual evapotranspiration $ET_a=(K_sK_{cb}+K_e)ET_0=5.69$ mm\,d$^{-1}$, a reduction of 26 per cent
\ans Relative yield by Doorenbos--Kassam with $K_y=1.25$: $1-Y/Y_m=K_y(1-ET_a/ET_c)=0.327$, i.e. a yield loss of 32.7 per cent if this stress lasts all season


}{}
\IfFileExists{extra/proj05.tex}{\banner{Design project 5: seasonal irrigation requirement of a maize field}
\textbf{Brief.} A maize crop is sown on 1 May on a loam. Monthly reference evapotranspiration and rainfall are given in Table~\ref{tab:nir5}. The task is to find the monthly crop evapotranspiration, the effective rainfall, the net irrigation requirement, the gross requirement for a sprinkler system (efficiency 75 per cent) and a drip system (90 per cent), the volume for a field of 30 ha, and the yield if the farmer supplies only 70 per cent of the net requirement.

\textbf{Method.} $K_c$ is evaluated at the middle day of each month; $ET_c=K_cET_0\times$days; effective rainfall is taken as 80 per cent of the monthly rain (capped at 60 mm); net irrigation requirement $NIR=\max(0,ET_c-P_{eff})$; gross $GIR=NIR/E_a$.

\begin{table}[!htbp]\centering\small\caption{Monthly computation of the net irrigation requirement of maize.}
\begin{tabular}{lrrrrrr}\toprule
\hd{Month} & \hd{$ET_0$ (mm/d)} & \hd{$K_c$} & \hd{$ET_c$ (mm)} & \hd{Rain (mm)} & \hd{$P_{eff}$ (mm)} & \hd{NIR (mm)}\\\midrule
May & 3.8 & 0.32 & 37 & 68 & 54 & 0\\
Jun & 4.9 & 1.20 & 176 & 52 & 42 & 135\\
Jul & 5.6 & 1.20 & 208 & 31 & 25 & 184\\
Aug & 5.9 & 0.60 & 110 & 24 & 19 & 91\\
Sep & 5.2 & 0.60 & 94 & 45 & 36 & 58\\
Oct & 4.1 & 0.60 & 76 & 64 & 51 & 25\\\bottomrule\end{tabular}\end{table}

\textbf{Season totals.} $ET_c=701$ mm, effective rain $=227$ mm, net requirement $NIR=492$ mm. A sprinkler system with $E_a=0.75$ needs $655$ mm gross, which over 30 ha is $196608$ m$^3$ ($197$ thousand m$^3$); a drip system with $E_a=0.90$ needs $546$ mm and $163840$ m$^3$, a saving of 17 per cent of the water of the sprinkler system.

\textbf{Deficit irrigation.} If only 70 per cent of $NIR$ is applied, $ET_a\approx0.7NIR+P_{eff}=571$ mm and $ET_a/ET_c=0.814$; by Equation~(\ref{eq:ky}) with $K_y=1.28$, the relative yield is $1-1.28\times(1-0.814)=0.762$, a loss of 23.8 per cent of a yield for a saving of 58982 m$^3$ of water on the 30 ha (sprinkler). The decision between full and deficit irrigation then depends on the price of water against the price of grain; the break-even price ratio is the loss in yield value per cubic metre saved.

}{}
\IfFileExists{extra/ten05.tex}{\subsection{Worked solutions on the ten soils}
The same season is irrigated and left rainfed on every soil, to show how the soil decides the gain.

\exbox{Ten-soil example 5a: Coarse loam (S01)}{Coarse loam: with $K_y=1.28$, $Y_m=11.5$ Mg\,ha$^{-1}$, compare the irrigated and rainfed seasons (125 d, $TAW=90$ mm) and give the yield, the water productivity and the irrigation water productivity.}\sol \givenline{$K_y=1.28$ ; $Y_m=11.5$ Mg\,ha$^{-1}$ ; $TAW=90$ mm}\ans Irrigated: $ET_a/ET_c=747/747=1.000$, $Y=11.5\,[1-1.28(1-1.000)]=11.50$ Mg\,ha$^{-1}$ \ans Rainfed: $ET_a/ET_c=0.321$, $Y=1.51$ \ans $WP=Y/ET_a=0.15$ kg\,m$^{-3}$ (irrigated) \ans Gain $9.99$ Mg\,ha$^{-1}$ for 577 mm of irrigation: $IWP=0.17$ kg\,m$^{-3}$ \finans

\exbox{Ten-soil example 5a: Coarse sand (S02)}{Coarse sand: with $K_y=1.28$, $Y_m=11.5$ Mg\,ha$^{-1}$, compare the irrigated and rainfed seasons (125 d, $TAW=30$ mm) and give the yield, the water productivity and the irrigation water productivity.}\sol \givenline{$K_y=1.28$ ; $Y_m=11.5$ Mg\,ha$^{-1}$ ; $TAW=30$ mm}\ans Irrigated: $ET_a/ET_c=747/747=1.000$, $Y=11.5\,[1-1.28(1-1.000)]=11.50$ Mg\,ha$^{-1}$ \ans Rainfed: $ET_a/ET_c=0.255$, $Y=0.53$ \ans $WP=Y/ET_a=0.15$ kg\,m$^{-3}$ (irrigated) \ans Gain $10.97$ Mg\,ha$^{-1}$ for 595 mm of irrigation: $IWP=0.18$ kg\,m$^{-3}$ \finans

\exbox{Ten-soil example 5a: Sand (S03)}{Sand: with $K_y=1.28$, $Y_m=11.5$ Mg\,ha$^{-1}$, compare the irrigated and rainfed seasons (125 d, $TAW=36$ mm) and give the yield, the water productivity and the irrigation water productivity.}\sol \givenline{$K_y=1.28$ ; $Y_m=11.5$ Mg\,ha$^{-1}$ ; $TAW=36$ mm}\ans Irrigated: $ET_a/ET_c=747/747=1.000$, $Y=11.5\,[1-1.28(1-1.000)]=11.50$ Mg\,ha$^{-1}$ \ans Rainfed: $ET_a/ET_c=0.262$, $Y=0.63$ \ans $WP=Y/ET_a=0.15$ kg\,m$^{-3}$ (irrigated) \ans Gain $10.87$ Mg\,ha$^{-1}$ for 579 mm of irrigation: $IWP=0.19$ kg\,m$^{-3}$ \finans

\exbox{Ten-soil example 5a: Loamy sand (S04)}{Loamy sand: with $K_y=1.28$, $Y_m=11.5$ Mg\,ha$^{-1}$, compare the irrigated and rainfed seasons (125 d, $TAW=54$ mm) and give the yield, the water productivity and the irrigation water productivity.}\sol \givenline{$K_y=1.28$ ; $Y_m=11.5$ Mg\,ha$^{-1}$ ; $TAW=54$ mm}\ans Irrigated: $ET_a/ET_c=747/747=1.000$, $Y=11.5\,[1-1.28(1-1.000)]=11.50$ Mg\,ha$^{-1}$ \ans Rainfed: $ET_a/ET_c=0.281$, $Y=0.92$ \ans $WP=Y/ET_a=0.15$ kg\,m$^{-3}$ (irrigated) \ans Gain $10.58$ Mg\,ha$^{-1}$ for 592 mm of irrigation: $IWP=0.18$ kg\,m$^{-3}$ \finans

\exbox{Ten-soil example 5a: Sandy loam (S05)}{Sandy loam: with $K_y=1.28$, $Y_m=11.5$ Mg\,ha$^{-1}$, compare the irrigated and rainfed seasons (125 d, $TAW=72$ mm) and give the yield, the water productivity and the irrigation water productivity.}\sol \givenline{$K_y=1.28$ ; $Y_m=11.5$ Mg\,ha$^{-1}$ ; $TAW=72$ mm}\ans Irrigated: $ET_a/ET_c=747/747=1.000$, $Y=11.5\,[1-1.28(1-1.000)]=11.50$ Mg\,ha$^{-1}$ \ans Rainfed: $ET_a/ET_c=0.301$, $Y=1.21$ \ans $WP=Y/ET_a=0.15$ kg\,m$^{-3}$ (irrigated) \ans Gain $10.29$ Mg\,ha$^{-1}$ for 557 mm of irrigation: $IWP=0.18$ kg\,m$^{-3}$ \finans

\exbox{Ten-soil example 5a: Loam (S06)}{Loam: with $K_y=1.28$, $Y_m=11.5$ Mg\,ha$^{-1}$, compare the irrigated and rainfed seasons (125 d, $TAW=102$ mm) and give the yield, the water productivity and the irrigation water productivity.}\sol \givenline{$K_y=1.28$ ; $Y_m=11.5$ Mg\,ha$^{-1}$ ; $TAW=102$ mm}\ans Irrigated: $ET_a/ET_c=747/747=1.000$, $Y=11.5\,[1-1.28(1-1.000)]=11.50$ Mg\,ha$^{-1}$ \ans Rainfed: $ET_a/ET_c=0.335$, $Y=1.71$ \ans $WP=Y/ET_a=0.15$ kg\,m$^{-3}$ (irrigated) \ans Gain $9.79$ Mg\,ha$^{-1}$ for 608 mm of irrigation: $IWP=0.16$ kg\,m$^{-3}$ \finans

\exbox{Ten-soil example 5a: Silt loam (S07)}{Silt loam: with $K_y=1.28$, $Y_m=11.5$ Mg\,ha$^{-1}$, compare the irrigated and rainfed seasons (125 d, $TAW=120$ mm) and give the yield, the water productivity and the irrigation water productivity.}\sol \givenline{$K_y=1.28$ ; $Y_m=11.5$ Mg\,ha$^{-1}$ ; $TAW=120$ mm}\ans Irrigated: $ET_a/ET_c=747/747=1.000$, $Y=11.5\,[1-1.28(1-1.000)]=11.50$ Mg\,ha$^{-1}$ \ans Rainfed: $ET_a/ET_c=0.356$, $Y=2.02$ \ans $WP=Y/ET_a=0.15$ kg\,m$^{-3}$ (irrigated) \ans Gain $9.48$ Mg\,ha$^{-1}$ for 573 mm of irrigation: $IWP=0.17$ kg\,m$^{-3}$ \finans

\exbox{Ten-soil example 5a: Sandy clay loam (S08)}{Sandy clay loam: with $K_y=1.28$, $Y_m=11.5$ Mg\,ha$^{-1}$, compare the irrigated and rainfed seasons (125 d, $TAW=84$ mm) and give the yield, the water productivity and the irrigation water productivity.}\sol \givenline{$K_y=1.28$ ; $Y_m=11.5$ Mg\,ha$^{-1}$ ; $TAW=84$ mm}\ans Irrigated: $ET_a/ET_c=747/747=1.000$, $Y=11.5\,[1-1.28(1-1.000)]=11.50$ Mg\,ha$^{-1}$ \ans Rainfed: $ET_a/ET_c=0.315$, $Y=1.41$ \ans $WP=Y/ET_a=0.15$ kg\,m$^{-3}$ (irrigated) \ans Gain $10.09$ Mg\,ha$^{-1}$ for 562 mm of irrigation: $IWP=0.18$ kg\,m$^{-3}$ \finans

\exbox{Ten-soil example 5a: Clay loam (S09)}{Clay loam: with $K_y=1.28$, $Y_m=11.5$ Mg\,ha$^{-1}$, compare the irrigated and rainfed seasons (125 d, $TAW=96$ mm) and give the yield, the water productivity and the irrigation water productivity.}\sol \givenline{$K_y=1.28$ ; $Y_m=11.5$ Mg\,ha$^{-1}$ ; $TAW=96$ mm}\ans Irrigated: $ET_a/ET_c=747/747=1.000$, $Y=11.5\,[1-1.28(1-1.000)]=11.50$ Mg\,ha$^{-1}$ \ans Rainfed: $ET_a/ET_c=0.328$, $Y=1.61$ \ans $WP=Y/ET_a=0.15$ kg\,m$^{-3}$ (irrigated) \ans Gain $9.89$ Mg\,ha$^{-1}$ for 561 mm of irrigation: $IWP=0.18$ kg\,m$^{-3}$ \finans

\exbox{Ten-soil example 5a: Silty clay loam (S10)}{Silty clay loam: with $K_y=1.28$, $Y_m=11.5$ Mg\,ha$^{-1}$, compare the irrigated and rainfed seasons (125 d, $TAW=102$ mm) and give the yield, the water productivity and the irrigation water productivity.}\sol \givenline{$K_y=1.28$ ; $Y_m=11.5$ Mg\,ha$^{-1}$ ; $TAW=102$ mm}\ans Irrigated: $ET_a/ET_c=747/747=1.000$, $Y=11.5\,[1-1.28(1-1.000)]=11.50$ Mg\,ha$^{-1}$ \ans Rainfed: $ET_a/ET_c=0.335$, $Y=1.71$ \ans $WP=Y/ET_a=0.15$ kg\,m$^{-3}$ (irrigated) \ans Gain $9.79$ Mg\,ha$^{-1}$ for 608 mm of irrigation: $IWP=0.16$ kg\,m$^{-3}$ \finans

}{}
\IfFileExists{extra/tenb05.tex}{\subsection{Worked solutions on ten crops}
The ten crops of the library differ in their demand, their roots and their sensitivity. The same type of question is answered for each of them at $ET_0=5.5$ mm\,d$^{-1}$; the data are typical values of the crops (Doorenbos--Kassam and FAO 56).

\exbox{Crop example 5a: Maize}{Maize ($K_y=1.25$, $Y_m=11.5$ Mg\,ha$^{-1}$): a deficit of 20 per cent of the seasonal evapotranspiration occurs ($\Sigma ET_m\approx660$ mm). Find the relative and absolute yield, the loss, and compare the value of the water saved (0.12 per m$^3$) with the grain lost (250 per Mg).}\sol $Y/Y_m=1-K_y(1-ET_a/ET_m)=1-1.25\times0.20=0.750$ \ans $Y=0.750\times11.5=8.62$ Mg\,ha$^{-1}$ \ans Loss $=2.88$ Mg\,ha$^{-1}$, worth 719 per ha \ans Water saved $=0.2\times6600=1320$ m$^3$\,ha$^{-1}$, worth 158 per ha \ans The grain lost is worth more than the water saved: the deficit does not pay at these prices.

\exbox{Crop example 5a: Wheat}{Wheat ($K_y=1.15$, $Y_m=8.0$ Mg\,ha$^{-1}$): a deficit of 20 per cent of the seasonal evapotranspiration occurs ($\Sigma ET_m\approx759$ mm). Find the relative and absolute yield, the loss, and compare the value of the water saved (0.12 per m$^3$) with the grain lost (250 per Mg).}\sol $Y/Y_m=1-K_y(1-ET_a/ET_m)=1-1.15\times0.20=0.770$ \ans $Y=0.770\times8.0=6.16$ Mg\,ha$^{-1}$ \ans Loss $=1.84$ Mg\,ha$^{-1}$, worth 460 per ha \ans Water saved $=0.2\times7590=1518$ m$^3$\,ha$^{-1}$, worth 182 per ha \ans The grain lost is worth more than the water saved: the deficit does not pay at these prices.

\exbox{Crop example 5a: Cotton}{Cotton ($K_y=0.85$, $Y_m=4.5$ Mg\,ha$^{-1}$): a deficit of 20 per cent of the seasonal evapotranspiration occurs ($\Sigma ET_m\approx911$ mm). Find the relative and absolute yield, the loss, and compare the value of the water saved (0.12 per m$^3$) with the grain lost (250 per Mg).}\sol $Y/Y_m=1-K_y(1-ET_a/ET_m)=1-0.85\times0.20=0.830$ \ans $Y=0.830\times4.5=3.73$ Mg\,ha$^{-1}$ \ans Loss $=0.77$ Mg\,ha$^{-1}$, worth 191 per ha \ans Water saved $=0.2\times9108=1822$ m$^3$\,ha$^{-1}$, worth 219 per ha \ans The saving pays.

\exbox{Crop example 5a: Tomato}{Tomato ($K_y=1.05$, $Y_m=75.0$ Mg\,ha$^{-1}$): a deficit of 20 per cent of the seasonal evapotranspiration occurs ($\Sigma ET_m\approx683$ mm). Find the relative and absolute yield, the loss, and compare the value of the water saved (0.12 per m$^3$) with the grain lost (250 per Mg).}\sol $Y/Y_m=1-K_y(1-ET_a/ET_m)=1-1.05\times0.20=0.790$ \ans $Y=0.790\times75.0=59.25$ Mg\,ha$^{-1}$ \ans Loss $=15.75$ Mg\,ha$^{-1}$, worth 3938 per ha \ans Water saved $=0.2\times6831=1366$ m$^3$\,ha$^{-1}$, worth 164 per ha \ans The grain lost is worth more than the water saved: the deficit does not pay at these prices.

\exbox{Crop example 5a: Potato}{Potato ($K_y=1.10$, $Y_m=40.0$ Mg\,ha$^{-1}$): a deficit of 20 per cent of the seasonal evapotranspiration occurs ($\Sigma ET_m\approx658$ mm). Find the relative and absolute yield, the loss, and compare the value of the water saved (0.12 per m$^3$) with the grain lost (250 per Mg).}\sol $Y/Y_m=1-K_y(1-ET_a/ET_m)=1-1.10\times0.20=0.780$ \ans $Y=0.780\times40.0=31.20$ Mg\,ha$^{-1}$ \ans Loss $=8.80$ Mg\,ha$^{-1}$, worth 2200 per ha \ans Water saved $=0.2\times6578=1316$ m$^3$\,ha$^{-1}$, worth 158 per ha \ans The grain lost is worth more than the water saved: the deficit does not pay at these prices.

\exbox{Crop example 5a: Soybean}{Soybean ($K_y=0.85$, $Y_m=3.5$ Mg\,ha$^{-1}$): a deficit of 20 per cent of the seasonal evapotranspiration occurs ($\Sigma ET_m\approx658$ mm). Find the relative and absolute yield, the loss, and compare the value of the water saved (0.12 per m$^3$) with the grain lost (250 per Mg).}\sol $Y/Y_m=1-K_y(1-ET_a/ET_m)=1-0.85\times0.20=0.830$ \ans $Y=0.830\times3.5=2.90$ Mg\,ha$^{-1}$ \ans Loss $=0.60$ Mg\,ha$^{-1}$, worth 149 per ha \ans Water saved $=0.2\times6578=1316$ m$^3$\,ha$^{-1}$, worth 158 per ha \ans The saving pays.

\exbox{Crop example 5a: Sugar beet}{Sugar beet ($K_y=1.00$, $Y_m=60.0$ Mg\,ha$^{-1}$): a deficit of 20 per cent of the seasonal evapotranspiration occurs ($\Sigma ET_m\approx950$ mm). Find the relative and absolute yield, the loss, and compare the value of the water saved (0.12 per m$^3$) with the grain lost (250 per Mg).}\sol $Y/Y_m=1-K_y(1-ET_a/ET_m)=1-1.00\times0.20=0.800$ \ans $Y=0.800\times60.0=48.00$ Mg\,ha$^{-1}$ \ans Loss $=12.00$ Mg\,ha$^{-1}$, worth 3000 per ha \ans Water saved $=0.2\times9504=1901$ m$^3$\,ha$^{-1}$, worth 228 per ha \ans The grain lost is worth more than the water saved: the deficit does not pay at these prices.

\exbox{Crop example 5a: Sunflower}{Sunflower ($K_y=0.95$, $Y_m=3.0$ Mg\,ha$^{-1}$): a deficit of 20 per cent of the seasonal evapotranspiration occurs ($\Sigma ET_m\approx605$ mm). Find the relative and absolute yield, the loss, and compare the value of the water saved (0.12 per m$^3$) with the grain lost (250 per Mg).}\sol $Y/Y_m=1-K_y(1-ET_a/ET_m)=1-0.95\times0.20=0.810$ \ans $Y=0.810\times3.0=2.43$ Mg\,ha$^{-1}$ \ans Loss $=0.57$ Mg\,ha$^{-1}$, worth 142 per ha \ans Water saved $=0.2\times6050=1210$ m$^3$\,ha$^{-1}$, worth 145 per ha \ans The saving pays.

\exbox{Crop example 5a: Onion}{Onion ($K_y=1.10$, $Y_m=50.0$ Mg\,ha$^{-1}$): a deficit of 20 per cent of the seasonal evapotranspiration occurs ($\Sigma ET_m\approx693$ mm). Find the relative and absolute yield, the loss, and compare the value of the water saved (0.12 per m$^3$) with the grain lost (250 per Mg).}\sol $Y/Y_m=1-K_y(1-ET_a/ET_m)=1-1.10\times0.20=0.780$ \ans $Y=0.780\times50.0=39.00$ Mg\,ha$^{-1}$ \ans Loss $=11.00$ Mg\,ha$^{-1}$, worth 2750 per ha \ans Water saved $=0.2\times6930=1386$ m$^3$\,ha$^{-1}$, worth 166 per ha \ans The grain lost is worth more than the water saved: the deficit does not pay at these prices.

\exbox{Crop example 5a: Alfalfa}{Alfalfa ($K_y=1.10$, $Y_m=15.0$ Mg\,ha$^{-1}$): a deficit of 20 per cent of the seasonal evapotranspiration occurs ($\Sigma ET_m\approx506$ mm). Find the relative and absolute yield, the loss, and compare the value of the water saved (0.12 per m$^3$) with the grain lost (250 per Mg).}\sol $Y/Y_m=1-K_y(1-ET_a/ET_m)=1-1.10\times0.20=0.780$ \ans $Y=0.780\times15.0=11.70$ Mg\,ha$^{-1}$ \ans Loss $=3.30$ Mg\,ha$^{-1}$, worth 825 per ha \ans Water saved $=0.2\times5060=1012$ m$^3$\,ha$^{-1}$, worth 121 per ha \ans The grain lost is worth more than the water saved: the deficit does not pay at these prices.

}{}
\IfFileExists{extra/teny05.tex}{\subsection{Worked solutions on the ten soils (third scenario)}
The same questions are repeated for different working conditions, to show how the answer moves with them.

\exbox{Ten-soil example 5c: Coarse loam (S01)}{Coarse loam: with $K_y=1.28$, $Y_m=11.5$ Mg\,ha$^{-1}$, compare the irrigated and rainfed seasons (125 d, $TAW=90$ mm) and give the yield, the water productivity and the irrigation water productivity.}\sol \givenline{$K_y=1.28$ ; $Y_m=11.5$ Mg\,ha$^{-1}$ ; $TAW=90$ mm}\ans Irrigated: $ET_a/ET_c=711/749=0.950$, \[Y=11.5\,[1-1.28(1-0.950)]=10.76\ \text{Mg\,ha$^{-1}$}\] \ans Rainfed: $ET_a/ET_c=0.279$, $Y=0.88$ \ans $WP=Y/ET_a=0.15$ kg\,m$^{-3}$ (irrigated) \ans Gain $9.88$ Mg\,ha$^{-1}$ for 582 mm of irrigation: $IWP=0.17$ kg\,m$^{-3}$ \finans

\exbox{Ten-soil example 5c: Coarse sand (S02)}{Coarse sand: with $K_y=1.28$, $Y_m=11.5$ Mg\,ha$^{-1}$, compare the irrigated and rainfed seasons (125 d, $TAW=30$ mm) and give the yield, the water productivity and the irrigation water productivity.}\sol \givenline{$K_y=1.28$ ; $Y_m=11.5$ Mg\,ha$^{-1}$ ; $TAW=30$ mm}\ans Irrigated: $ET_a/ET_c=716/749=0.956$, \[Y=11.5\,[1-1.28(1-0.956)]=10.86\ \text{Mg\,ha$^{-1}$}\] \ans Rainfed: $ET_a/ET_c=0.207$, $Y=-0.17$ \ans $WP=Y/ET_a=0.15$ kg\,m$^{-3}$ (irrigated) \ans Gain $11.03$ Mg\,ha$^{-1}$ for 592 mm of irrigation: $IWP=0.19$ kg\,m$^{-3}$ \finans

\exbox{Ten-soil example 5c: Sand (S03)}{Sand: with $K_y=1.28$, $Y_m=11.5$ Mg\,ha$^{-1}$, compare the irrigated and rainfed seasons (125 d, $TAW=36$ mm) and give the yield, the water productivity and the irrigation water productivity.}\sol \givenline{$K_y=1.28$ ; $Y_m=11.5$ Mg\,ha$^{-1}$ ; $TAW=36$ mm}\ans Irrigated: $ET_a/ET_c=712/749=0.951$, \[Y=11.5\,[1-1.28(1-0.951)]=10.78\ \text{Mg\,ha$^{-1}$}\] \ans Rainfed: $ET_a/ET_c=0.218$, $Y=-0.00$ \ans $WP=Y/ET_a=0.15$ kg\,m$^{-3}$ (irrigated) \ans Gain $10.78$ Mg\,ha$^{-1}$ for 612 mm of irrigation: $IWP=0.18$ kg\,m$^{-3}$ \finans

\exbox{Ten-soil example 5c: Loamy sand (S04)}{Loamy sand: with $K_y=1.28$, $Y_m=11.5$ Mg\,ha$^{-1}$, compare the irrigated and rainfed seasons (125 d, $TAW=54$ mm) and give the yield, the water productivity and the irrigation water productivity.}\sol \givenline{$K_y=1.28$ ; $Y_m=11.5$ Mg\,ha$^{-1}$ ; $TAW=54$ mm}\ans Irrigated: $ET_a/ET_c=714/749=0.954$, \[Y=11.5\,[1-1.28(1-0.954)]=10.82\ \text{Mg\,ha$^{-1}$}\] \ans Rainfed: $ET_a/ET_c=0.239$, $Y=0.30$ \ans $WP=Y/ET_a=0.15$ kg\,m$^{-3}$ (irrigated) \ans Gain $10.52$ Mg\,ha$^{-1}$ for 596 mm of irrigation: $IWP=0.18$ kg\,m$^{-3}$ \finans

\exbox{Ten-soil example 5c: Sandy loam (S05)}{Sandy loam: with $K_y=1.28$, $Y_m=11.5$ Mg\,ha$^{-1}$, compare the irrigated and rainfed seasons (125 d, $TAW=72$ mm) and give the yield, the water productivity and the irrigation water productivity.}\sol \givenline{$K_y=1.28$ ; $Y_m=11.5$ Mg\,ha$^{-1}$ ; $TAW=72$ mm}\ans Irrigated: $ET_a/ET_c=701/749=0.936$, \[Y=11.5\,[1-1.28(1-0.936)]=10.56\ \text{Mg\,ha$^{-1}$}\] \ans Rainfed: $ET_a/ET_c=0.259$, $Y=0.60$ \ans $WP=Y/ET_a=0.15$ kg\,m$^{-3}$ (irrigated) \ans Gain $9.96$ Mg\,ha$^{-1}$ for 587 mm of irrigation: $IWP=0.17$ kg\,m$^{-3}$ \finans

\exbox{Ten-soil example 5c: Loam (S06)}{Loam: with $K_y=1.28$, $Y_m=11.5$ Mg\,ha$^{-1}$, compare the irrigated and rainfed seasons (125 d, $TAW=102$ mm) and give the yield, the water productivity and the irrigation water productivity.}\sol \givenline{$K_y=1.28$ ; $Y_m=11.5$ Mg\,ha$^{-1}$ ; $TAW=102$ mm}\ans Irrigated: $ET_a/ET_c=701/749=0.937$, \[Y=11.5\,[1-1.28(1-0.937)]=10.57\ \text{Mg\,ha$^{-1}$}\] \ans Rainfed: $ET_a/ET_c=0.291$, $Y=1.07$ \ans $WP=Y/ET_a=0.15$ kg\,m$^{-3}$ (irrigated) \ans Gain $9.50$ Mg\,ha$^{-1}$ for 519 mm of irrigation: $IWP=0.18$ kg\,m$^{-3}$ \finans

\exbox{Ten-soil example 5c: Silt loam (S07)}{Silt loam: with $K_y=1.28$, $Y_m=11.5$ Mg\,ha$^{-1}$, compare the irrigated and rainfed seasons (125 d, $TAW=120$ mm) and give the yield, the water productivity and the irrigation water productivity.}\sol \givenline{$K_y=1.28$ ; $Y_m=11.5$ Mg\,ha$^{-1}$ ; $TAW=120$ mm}\ans Irrigated: $ET_a/ET_c=705/749=0.941$, \[Y=11.5\,[1-1.28(1-0.941)]=10.64\ \text{Mg\,ha$^{-1}$}\] \ans Rainfed: $ET_a/ET_c=0.311$, $Y=1.36$ \ans $WP=Y/ET_a=0.15$ kg\,m$^{-3}$ (irrigated) \ans Gain $9.28$ Mg\,ha$^{-1}$ for 520 mm of irrigation: $IWP=0.18$ kg\,m$^{-3}$ \finans

\exbox{Ten-soil example 5c: Sandy clay loam (S08)}{Sandy clay loam: with $K_y=1.28$, $Y_m=11.5$ Mg\,ha$^{-1}$, compare the irrigated and rainfed seasons (125 d, $TAW=84$ mm) and give the yield, the water productivity and the irrigation water productivity.}\sol \givenline{$K_y=1.28$ ; $Y_m=11.5$ Mg\,ha$^{-1}$ ; $TAW=84$ mm}\ans Irrigated: $ET_a/ET_c=705/749=0.941$, \[Y=11.5\,[1-1.28(1-0.941)]=10.63\ \text{Mg\,ha$^{-1}$}\] \ans Rainfed: $ET_a/ET_c=0.272$, $Y=0.79$ \ans $WP=Y/ET_a=0.15$ kg\,m$^{-3}$ (irrigated) \ans Gain $9.85$ Mg\,ha$^{-1}$ for 552 mm of irrigation: $IWP=0.18$ kg\,m$^{-3}$ \finans

\exbox{Ten-soil example 5c: Clay loam (S09)}{Clay loam: with $K_y=1.28$, $Y_m=11.5$ Mg\,ha$^{-1}$, compare the irrigated and rainfed seasons (125 d, $TAW=96$ mm) and give the yield, the water productivity and the irrigation water productivity.}\sol \givenline{$K_y=1.28$ ; $Y_m=11.5$ Mg\,ha$^{-1}$ ; $TAW=96$ mm}\ans Irrigated: $ET_a/ET_c=705/749=0.942$, \[Y=11.5\,[1-1.28(1-0.942)]=10.64\ \text{Mg\,ha$^{-1}$}\] \ans Rainfed: $ET_a/ET_c=0.285$, $Y=0.98$ \ans $WP=Y/ET_a=0.15$ kg\,m$^{-3}$ (irrigated) \ans Gain $9.67$ Mg\,ha$^{-1}$ for 554 mm of irrigation: $IWP=0.17$ kg\,m$^{-3}$ \finans

\exbox{Ten-soil example 5c: Silty clay loam (S10)}{Silty clay loam: with $K_y=1.28$, $Y_m=11.5$ Mg\,ha$^{-1}$, compare the irrigated and rainfed seasons (125 d, $TAW=102$ mm) and give the yield, the water productivity and the irrigation water productivity.}\sol \givenline{$K_y=1.28$ ; $Y_m=11.5$ Mg\,ha$^{-1}$ ; $TAW=102$ mm}\ans Irrigated: $ET_a/ET_c=701/749=0.937$, \[Y=11.5\,[1-1.28(1-0.937)]=10.57\ \text{Mg\,ha$^{-1}$}\] \ans Rainfed: $ET_a/ET_c=0.291$, $Y=1.07$ \ans $WP=Y/ET_a=0.15$ kg\,m$^{-3}$ (irrigated) \ans Gain $9.50$ Mg\,ha$^{-1}$ for 519 mm of irrigation: $IWP=0.18$ kg\,m$^{-3}$ \finans

}{}
\IfFileExists{extra/tenby05.tex}{\section{Worked solutions on ten crops (hot climate)}
The ten crops of the library differ in their demand, their roots and their sensitivity. The same type of question is answered for each of them at $ET_0=7.5$ mm\,d$^{-1}$; the data are typical values of the crops (Doorenbos--Kassam and FAO 56).

\exbox{Crop example 5c: Maize}{Maize ($K_y=1.25$, $Y_m=11.5$ Mg\,ha$^{-1}$): a deficit of 20 per cent of the seasonal evapotranspiration occurs ($\Sigma ET_m\approx900$ mm). Find the relative and absolute yield, the loss, and compare the value of the water saved (0.12 per m$^3$) with the grain lost (250 per Mg).}\sol $Y/Y_m=1-K_y(1-ET_a/ET_m)=1-1.25\times0.20=0.750$ \ans $Y=0.750\times11.5=8.62$ Mg\,ha$^{-1}$ \ans Loss $=2.88$ Mg\,ha$^{-1}$, worth 719 per ha \ans Water saved $=0.2\times9000=1800$ m$^3$\,ha$^{-1}$, worth 216 per ha \ans The grain lost is worth more than the water saved: the deficit does not pay at these prices.

\exbox{Crop example 5c: Wheat}{Wheat ($K_y=1.15$, $Y_m=8.0$ Mg\,ha$^{-1}$): a deficit of 20 per cent of the seasonal evapotranspiration occurs ($\Sigma ET_m\approx1035$ mm). Find the relative and absolute yield, the loss, and compare the value of the water saved (0.12 per m$^3$) with the grain lost (250 per Mg).}\sol $Y/Y_m=1-K_y(1-ET_a/ET_m)=1-1.15\times0.20=0.770$ \ans $Y=0.770\times8.0=6.16$ Mg\,ha$^{-1}$ \ans Loss $=1.84$ Mg\,ha$^{-1}$, worth 460 per ha \ans Water saved $=0.2\times10350=2070$ m$^3$\,ha$^{-1}$, worth 248 per ha \ans The grain lost is worth more than the water saved: the deficit does not pay at these prices.

\exbox{Crop example 5c: Cotton}{Cotton ($K_y=0.85$, $Y_m=4.5$ Mg\,ha$^{-1}$): a deficit of 20 per cent of the seasonal evapotranspiration occurs ($\Sigma ET_m\approx1242$ mm). Find the relative and absolute yield, the loss, and compare the value of the water saved (0.12 per m$^3$) with the grain lost (250 per Mg).}\sol $Y/Y_m=1-K_y(1-ET_a/ET_m)=1-0.85\times0.20=0.830$ \ans $Y=0.830\times4.5=3.73$ Mg\,ha$^{-1}$ \ans Loss $=0.77$ Mg\,ha$^{-1}$, worth 191 per ha \ans Water saved $=0.2\times12420=2484$ m$^3$\,ha$^{-1}$, worth 298 per ha \ans The saving pays.

\exbox{Crop example 5c: Tomato}{Tomato ($K_y=1.05$, $Y_m=75.0$ Mg\,ha$^{-1}$): a deficit of 20 per cent of the seasonal evapotranspiration occurs ($\Sigma ET_m\approx932$ mm). Find the relative and absolute yield, the loss, and compare the value of the water saved (0.12 per m$^3$) with the grain lost (250 per Mg).}\sol $Y/Y_m=1-K_y(1-ET_a/ET_m)=1-1.05\times0.20=0.790$ \ans $Y=0.790\times75.0=59.25$ Mg\,ha$^{-1}$ \ans Loss $=15.75$ Mg\,ha$^{-1}$, worth 3938 per ha \ans Water saved $=0.2\times9315=1863$ m$^3$\,ha$^{-1}$, worth 224 per ha \ans The grain lost is worth more than the water saved: the deficit does not pay at these prices.

\exbox{Crop example 5c: Potato}{Potato ($K_y=1.10$, $Y_m=40.0$ Mg\,ha$^{-1}$): a deficit of 20 per cent of the seasonal evapotranspiration occurs ($\Sigma ET_m\approx897$ mm). Find the relative and absolute yield, the loss, and compare the value of the water saved (0.12 per m$^3$) with the grain lost (250 per Mg).}\sol $Y/Y_m=1-K_y(1-ET_a/ET_m)=1-1.10\times0.20=0.780$ \ans $Y=0.780\times40.0=31.20$ Mg\,ha$^{-1}$ \ans Loss $=8.80$ Mg\,ha$^{-1}$, worth 2200 per ha \ans Water saved $=0.2\times8970=1794$ m$^3$\,ha$^{-1}$, worth 215 per ha \ans The grain lost is worth more than the water saved: the deficit does not pay at these prices.

\exbox{Crop example 5c: Soybean}{Soybean ($K_y=0.85$, $Y_m=3.5$ Mg\,ha$^{-1}$): a deficit of 20 per cent of the seasonal evapotranspiration occurs ($\Sigma ET_m\approx897$ mm). Find the relative and absolute yield, the loss, and compare the value of the water saved (0.12 per m$^3$) with the grain lost (250 per Mg).}\sol $Y/Y_m=1-K_y(1-ET_a/ET_m)=1-0.85\times0.20=0.830$ \ans $Y=0.830\times3.5=2.90$ Mg\,ha$^{-1}$ \ans Loss $=0.60$ Mg\,ha$^{-1}$, worth 149 per ha \ans Water saved $=0.2\times8970=1794$ m$^3$\,ha$^{-1}$, worth 215 per ha \ans The saving pays.

\exbox{Crop example 5c: Sugar beet}{Sugar beet ($K_y=1.00$, $Y_m=60.0$ Mg\,ha$^{-1}$): a deficit of 20 per cent of the seasonal evapotranspiration occurs ($\Sigma ET_m\approx1296$ mm). Find the relative and absolute yield, the loss, and compare the value of the water saved (0.12 per m$^3$) with the grain lost (250 per Mg).}\sol $Y/Y_m=1-K_y(1-ET_a/ET_m)=1-1.00\times0.20=0.800$ \ans $Y=0.800\times60.0=48.00$ Mg\,ha$^{-1}$ \ans Loss $=12.00$ Mg\,ha$^{-1}$, worth 3000 per ha \ans Water saved $=0.2\times12960=2592$ m$^3$\,ha$^{-1}$, worth 311 per ha \ans The grain lost is worth more than the water saved: the deficit does not pay at these prices.

\exbox{Crop example 5c: Sunflower}{Sunflower ($K_y=0.95$, $Y_m=3.0$ Mg\,ha$^{-1}$): a deficit of 20 per cent of the seasonal evapotranspiration occurs ($\Sigma ET_m\approx825$ mm). Find the relative and absolute yield, the loss, and compare the value of the water saved (0.12 per m$^3$) with the grain lost (250 per Mg).}\sol $Y/Y_m=1-K_y(1-ET_a/ET_m)=1-0.95\times0.20=0.810$ \ans $Y=0.810\times3.0=2.43$ Mg\,ha$^{-1}$ \ans Loss $=0.57$ Mg\,ha$^{-1}$, worth 142 per ha \ans Water saved $=0.2\times8250=1650$ m$^3$\,ha$^{-1}$, worth 198 per ha \ans The saving pays.

\exbox{Crop example 5c: Onion}{Onion ($K_y=1.10$, $Y_m=50.0$ Mg\,ha$^{-1}$): a deficit of 20 per cent of the seasonal evapotranspiration occurs ($\Sigma ET_m\approx945$ mm). Find the relative and absolute yield, the loss, and compare the value of the water saved (0.12 per m$^3$) with the grain lost (250 per Mg).}\sol $Y/Y_m=1-K_y(1-ET_a/ET_m)=1-1.10\times0.20=0.780$ \ans $Y=0.780\times50.0=39.00$ Mg\,ha$^{-1}$ \ans Loss $=11.00$ Mg\,ha$^{-1}$, worth 2750 per ha \ans Water saved $=0.2\times9450=1890$ m$^3$\,ha$^{-1}$, worth 227 per ha \ans The grain lost is worth more than the water saved: the deficit does not pay at these prices.

\exbox{Crop example 5c: Alfalfa}{Alfalfa ($K_y=1.10$, $Y_m=15.0$ Mg\,ha$^{-1}$): a deficit of 20 per cent of the seasonal evapotranspiration occurs ($\Sigma ET_m\approx690$ mm). Find the relative and absolute yield, the loss, and compare the value of the water saved (0.12 per m$^3$) with the grain lost (250 per Mg).}\sol $Y/Y_m=1-K_y(1-ET_a/ET_m)=1-1.10\times0.20=0.780$ \ans $Y=0.780\times15.0=11.70$ Mg\,ha$^{-1}$ \ans Loss $=3.30$ Mg\,ha$^{-1}$, worth 825 per ha \ans Water saved $=0.2\times6900=1380$ m$^3$\,ha$^{-1}$, worth 166 per ha \ans The grain lost is worth more than the water saved: the deficit does not pay at these prices.

}{}
\IfFileExists{extra/tenx05.tex}{\subsection{Worked solutions on the ten soils (second scenario)}
The same questions are repeated for different working conditions, to show how the answer moves with them.

\exbox{Ten-soil example 5b: Coarse loam (S01)}{Coarse loam: with $K_y=1.28$, $Y_m=11.5$ Mg\,ha$^{-1}$, compare the irrigated and rainfed seasons (125 d, $TAW=90$ mm) and give the yield, the water productivity and the irrigation water productivity.}\sol Irrigated: $ET_a/ET_c=759/759=1.000$, $Y=11.5\,[1-1.28(1-1.000)]=11.50$ Mg\,ha$^{-1}$ \ans Rainfed: $ET_a/ET_c=0.320$, $Y=1.50$ \ans $WP=Y/ET_a=0.15$ kg\,m$^{-3}$ (irrigated) \ans Gain $10.00$ Mg\,ha$^{-1}$ for 569 mm of irrigation: $IWP=0.18$ kg\,m$^{-3}$ \ans

\exbox{Ten-soil example 5b: Coarse sand (S02)}{Coarse sand: with $K_y=1.28$, $Y_m=11.5$ Mg\,ha$^{-1}$, compare the irrigated and rainfed seasons (125 d, $TAW=30$ mm) and give the yield, the water productivity and the irrigation water productivity.}\sol Irrigated: $ET_a/ET_c=759/759=1.000$, $Y=11.5\,[1-1.28(1-1.000)]=11.50$ Mg\,ha$^{-1}$ \ans Rainfed: $ET_a/ET_c=0.262$, $Y=0.64$ \ans $WP=Y/ET_a=0.15$ kg\,m$^{-3}$ (irrigated) \ans Gain $10.86$ Mg\,ha$^{-1}$ for 615 mm of irrigation: $IWP=0.18$ kg\,m$^{-3}$ \ans

\exbox{Ten-soil example 5b: Sand (S03)}{Sand: with $K_y=1.28$, $Y_m=11.5$ Mg\,ha$^{-1}$, compare the irrigated and rainfed seasons (125 d, $TAW=36$ mm) and give the yield, the water productivity and the irrigation water productivity.}\sol Irrigated: $ET_a/ET_c=759/759=1.000$, $Y=11.5\,[1-1.28(1-1.000)]=11.50$ Mg\,ha$^{-1}$ \ans Rainfed: $ET_a/ET_c=0.268$, $Y=0.72$ \ans $WP=Y/ET_a=0.15$ kg\,m$^{-3}$ (irrigated) \ans Gain $10.78$ Mg\,ha$^{-1}$ for 596 mm of irrigation: $IWP=0.18$ kg\,m$^{-3}$ \ans

\exbox{Ten-soil example 5b: Loamy sand (S04)}{Loamy sand: with $K_y=1.28$, $Y_m=11.5$ Mg\,ha$^{-1}$, compare the irrigated and rainfed seasons (125 d, $TAW=54$ mm) and give the yield, the water productivity and the irrigation water productivity.}\sol Irrigated: $ET_a/ET_c=759/759=1.000$, $Y=11.5\,[1-1.28(1-1.000)]=11.50$ Mg\,ha$^{-1}$ \ans Rainfed: $ET_a/ET_c=0.284$, $Y=0.96$ \ans $WP=Y/ET_a=0.15$ kg\,m$^{-3}$ (irrigated) \ans Gain $10.54$ Mg\,ha$^{-1}$ for 589 mm of irrigation: $IWP=0.18$ kg\,m$^{-3}$ \ans

\exbox{Ten-soil example 5b: Sandy loam (S05)}{Sandy loam: with $K_y=1.28$, $Y_m=11.5$ Mg\,ha$^{-1}$, compare the irrigated and rainfed seasons (125 d, $TAW=72$ mm) and give the yield, the water productivity and the irrigation water productivity.}\sol Irrigated: $ET_a/ET_c=759/759=1.000$, $Y=11.5\,[1-1.28(1-1.000)]=11.50$ Mg\,ha$^{-1}$ \ans Rainfed: $ET_a/ET_c=0.302$, $Y=1.22$ \ans $WP=Y/ET_a=0.15$ kg\,m$^{-3}$ (irrigated) \ans Gain $10.28$ Mg\,ha$^{-1}$ for 599 mm of irrigation: $IWP=0.17$ kg\,m$^{-3}$ \ans

\exbox{Ten-soil example 5b: Loam (S06)}{Loam: with $K_y=1.28$, $Y_m=11.5$ Mg\,ha$^{-1}$, compare the irrigated and rainfed seasons (125 d, $TAW=102$ mm) and give the yield, the water productivity and the irrigation water productivity.}\sol Irrigated: $ET_a/ET_c=759/759=1.000$, $Y=11.5\,[1-1.28(1-1.000)]=11.50$ Mg\,ha$^{-1}$ \ans Rainfed: $ET_a/ET_c=0.333$, $Y=1.68$ \ans $WP=Y/ET_a=0.15$ kg\,m$^{-3}$ (irrigated) \ans Gain $9.82$ Mg\,ha$^{-1}$ for 574 mm of irrigation: $IWP=0.17$ kg\,m$^{-3}$ \ans

\exbox{Ten-soil example 5b: Silt loam (S07)}{Silt loam: with $K_y=1.28$, $Y_m=11.5$ Mg\,ha$^{-1}$, compare the irrigated and rainfed seasons (125 d, $TAW=120$ mm) and give the yield, the water productivity and the irrigation water productivity.}\sol Irrigated: $ET_a/ET_c=759/759=1.000$, $Y=11.5\,[1-1.28(1-1.000)]=11.50$ Mg\,ha$^{-1}$ \ans Rainfed: $ET_a/ET_c=0.353$, $Y=1.97$ \ans $WP=Y/ET_a=0.15$ kg\,m$^{-3}$ (irrigated) \ans Gain $9.53$ Mg\,ha$^{-1}$ for 558 mm of irrigation: $IWP=0.17$ kg\,m$^{-3}$ \ans

\exbox{Ten-soil example 5b: Sandy clay loam (S08)}{Sandy clay loam: with $K_y=1.28$, $Y_m=11.5$ Mg\,ha$^{-1}$, compare the irrigated and rainfed seasons (125 d, $TAW=84$ mm) and give the yield, the water productivity and the irrigation water productivity.}\sol Irrigated: $ET_a/ET_c=759/759=1.000$, $Y=11.5\,[1-1.28(1-1.000)]=11.50$ Mg\,ha$^{-1}$ \ans Rainfed: $ET_a/ET_c=0.314$, $Y=1.40$ \ans $WP=Y/ET_a=0.15$ kg\,m$^{-3}$ (irrigated) \ans Gain $10.10$ Mg\,ha$^{-1}$ for 570 mm of irrigation: $IWP=0.18$ kg\,m$^{-3}$ \ans

\exbox{Ten-soil example 5b: Clay loam (S09)}{Clay loam: with $K_y=1.28$, $Y_m=11.5$ Mg\,ha$^{-1}$, compare the irrigated and rainfed seasons (125 d, $TAW=96$ mm) and give the yield, the water productivity and the irrigation water productivity.}\sol Irrigated: $ET_a/ET_c=759/759=1.000$, $Y=11.5\,[1-1.28(1-1.000)]=11.50$ Mg\,ha$^{-1}$ \ans Rainfed: $ET_a/ET_c=0.327$, $Y=1.59$ \ans $WP=Y/ET_a=0.15$ kg\,m$^{-3}$ (irrigated) \ans Gain $9.91$ Mg\,ha$^{-1}$ for 582 mm of irrigation: $IWP=0.17$ kg\,m$^{-3}$ \ans

\exbox{Ten-soil example 5b: Silty clay loam (S10)}{Silty clay loam: with $K_y=1.28$, $Y_m=11.5$ Mg\,ha$^{-1}$, compare the irrigated and rainfed seasons (125 d, $TAW=102$ mm) and give the yield, the water productivity and the irrigation water productivity.}\sol Irrigated: $ET_a/ET_c=759/759=1.000$, $Y=11.5\,[1-1.28(1-1.000)]=11.50$ Mg\,ha$^{-1}$ \ans Rainfed: $ET_a/ET_c=0.333$, $Y=1.68$ \ans $WP=Y/ET_a=0.15$ kg\,m$^{-3}$ (irrigated) \ans Gain $9.82$ Mg\,ha$^{-1}$ for 574 mm of irrigation: $IWP=0.17$ kg\,m$^{-3}$ \ans

}{}
\IfFileExists{extra/tenbx05.tex}{\subsection{Worked solutions on ten crops (cool climate)}
The ten crops of the library differ in their demand, their roots and their sensitivity. The same type of question is answered for each of them at $ET_0=4.0$ mm\,d$^{-1}$; the data are typical values of the crops (Doorenbos--Kassam and FAO 56).

\exbox{Crop example 5b: Maize}{Maize ($K_y=1.25$, $Y_m=11.5$ Mg\,ha$^{-1}$): a deficit of 20 per cent of the seasonal evapotranspiration occurs ($\Sigma ET_m\approx480$ mm). Find the relative and absolute yield, the loss, and compare the value of the water saved (0.12 per m$^3$) with the grain lost (250 per Mg).}\sol $Y/Y_m=1-K_y(1-ET_a/ET_m)=1-1.25\times0.20=0.750$ \ans $Y=0.750\times11.5=8.62$ Mg\,ha$^{-1}$ \ans Loss $=2.88$ Mg\,ha$^{-1}$, worth 719 per ha \ans Water saved $=0.2\times4800=960$ m$^3$\,ha$^{-1}$, worth 115 per ha \ans The grain lost is worth more than the water saved: the deficit does not pay at these prices.

\exbox{Crop example 5b: Wheat}{Wheat ($K_y=1.15$, $Y_m=8.0$ Mg\,ha$^{-1}$): a deficit of 20 per cent of the seasonal evapotranspiration occurs ($\Sigma ET_m\approx552$ mm). Find the relative and absolute yield, the loss, and compare the value of the water saved (0.12 per m$^3$) with the grain lost (250 per Mg).}\sol $Y/Y_m=1-K_y(1-ET_a/ET_m)=1-1.15\times0.20=0.770$ \ans $Y=0.770\times8.0=6.16$ Mg\,ha$^{-1}$ \ans Loss $=1.84$ Mg\,ha$^{-1}$, worth 460 per ha \ans Water saved $=0.2\times5520=1104$ m$^3$\,ha$^{-1}$, worth 132 per ha \ans The grain lost is worth more than the water saved: the deficit does not pay at these prices.

\exbox{Crop example 5b: Cotton}{Cotton ($K_y=0.85$, $Y_m=4.5$ Mg\,ha$^{-1}$): a deficit of 20 per cent of the seasonal evapotranspiration occurs ($\Sigma ET_m\approx662$ mm). Find the relative and absolute yield, the loss, and compare the value of the water saved (0.12 per m$^3$) with the grain lost (250 per Mg).}\sol $Y/Y_m=1-K_y(1-ET_a/ET_m)=1-0.85\times0.20=0.830$ \ans $Y=0.830\times4.5=3.73$ Mg\,ha$^{-1}$ \ans Loss $=0.77$ Mg\,ha$^{-1}$, worth 191 per ha \ans Water saved $=0.2\times6624=1325$ m$^3$\,ha$^{-1}$, worth 159 per ha \ans The grain lost is worth more than the water saved: the deficit does not pay at these prices.

\exbox{Crop example 5b: Tomato}{Tomato ($K_y=1.05$, $Y_m=75.0$ Mg\,ha$^{-1}$): a deficit of 20 per cent of the seasonal evapotranspiration occurs ($\Sigma ET_m\approx497$ mm). Find the relative and absolute yield, the loss, and compare the value of the water saved (0.12 per m$^3$) with the grain lost (250 per Mg).}\sol $Y/Y_m=1-K_y(1-ET_a/ET_m)=1-1.05\times0.20=0.790$ \ans $Y=0.790\times75.0=59.25$ Mg\,ha$^{-1}$ \ans Loss $=15.75$ Mg\,ha$^{-1}$, worth 3938 per ha \ans Water saved $=0.2\times4968=994$ m$^3$\,ha$^{-1}$, worth 119 per ha \ans The grain lost is worth more than the water saved: the deficit does not pay at these prices.

\exbox{Crop example 5b: Potato}{Potato ($K_y=1.10$, $Y_m=40.0$ Mg\,ha$^{-1}$): a deficit of 20 per cent of the seasonal evapotranspiration occurs ($\Sigma ET_m\approx478$ mm). Find the relative and absolute yield, the loss, and compare the value of the water saved (0.12 per m$^3$) with the grain lost (250 per Mg).}\sol $Y/Y_m=1-K_y(1-ET_a/ET_m)=1-1.10\times0.20=0.780$ \ans $Y=0.780\times40.0=31.20$ Mg\,ha$^{-1}$ \ans Loss $=8.80$ Mg\,ha$^{-1}$, worth 2200 per ha \ans Water saved $=0.2\times4784=957$ m$^3$\,ha$^{-1}$, worth 115 per ha \ans The grain lost is worth more than the water saved: the deficit does not pay at these prices.

\exbox{Crop example 5b: Soybean}{Soybean ($K_y=0.85$, $Y_m=3.5$ Mg\,ha$^{-1}$): a deficit of 20 per cent of the seasonal evapotranspiration occurs ($\Sigma ET_m\approx478$ mm). Find the relative and absolute yield, the loss, and compare the value of the water saved (0.12 per m$^3$) with the grain lost (250 per Mg).}\sol $Y/Y_m=1-K_y(1-ET_a/ET_m)=1-0.85\times0.20=0.830$ \ans $Y=0.830\times3.5=2.90$ Mg\,ha$^{-1}$ \ans Loss $=0.60$ Mg\,ha$^{-1}$, worth 149 per ha \ans Water saved $=0.2\times4784=957$ m$^3$\,ha$^{-1}$, worth 115 per ha \ans The grain lost is worth more than the water saved: the deficit does not pay at these prices.

\exbox{Crop example 5b: Sugar beet}{Sugar beet ($K_y=1.00$, $Y_m=60.0$ Mg\,ha$^{-1}$): a deficit of 20 per cent of the seasonal evapotranspiration occurs ($\Sigma ET_m\approx691$ mm). Find the relative and absolute yield, the loss, and compare the value of the water saved (0.12 per m$^3$) with the grain lost (250 per Mg).}\sol $Y/Y_m=1-K_y(1-ET_a/ET_m)=1-1.00\times0.20=0.800$ \ans $Y=0.800\times60.0=48.00$ Mg\,ha$^{-1}$ \ans Loss $=12.00$ Mg\,ha$^{-1}$, worth 3000 per ha \ans Water saved $=0.2\times6912=1382$ m$^3$\,ha$^{-1}$, worth 166 per ha \ans The grain lost is worth more than the water saved: the deficit does not pay at these prices.

\exbox{Crop example 5b: Sunflower}{Sunflower ($K_y=0.95$, $Y_m=3.0$ Mg\,ha$^{-1}$): a deficit of 20 per cent of the seasonal evapotranspiration occurs ($\Sigma ET_m\approx440$ mm). Find the relative and absolute yield, the loss, and compare the value of the water saved (0.12 per m$^3$) with the grain lost (250 per Mg).}\sol $Y/Y_m=1-K_y(1-ET_a/ET_m)=1-0.95\times0.20=0.810$ \ans $Y=0.810\times3.0=2.43$ Mg\,ha$^{-1}$ \ans Loss $=0.57$ Mg\,ha$^{-1}$, worth 142 per ha \ans Water saved $=0.2\times4400=880$ m$^3$\,ha$^{-1}$, worth 106 per ha \ans The grain lost is worth more than the water saved: the deficit does not pay at these prices.

\exbox{Crop example 5b: Onion}{Onion ($K_y=1.10$, $Y_m=50.0$ Mg\,ha$^{-1}$): a deficit of 20 per cent of the seasonal evapotranspiration occurs ($\Sigma ET_m\approx504$ mm). Find the relative and absolute yield, the loss, and compare the value of the water saved (0.12 per m$^3$) with the grain lost (250 per Mg).}\sol $Y/Y_m=1-K_y(1-ET_a/ET_m)=1-1.10\times0.20=0.780$ \ans $Y=0.780\times50.0=39.00$ Mg\,ha$^{-1}$ \ans Loss $=11.00$ Mg\,ha$^{-1}$, worth 2750 per ha \ans Water saved $=0.2\times5040=1008$ m$^3$\,ha$^{-1}$, worth 121 per ha \ans The grain lost is worth more than the water saved: the deficit does not pay at these prices.

\exbox{Crop example 5b: Alfalfa}{Alfalfa ($K_y=1.10$, $Y_m=15.0$ Mg\,ha$^{-1}$): a deficit of 20 per cent of the seasonal evapotranspiration occurs ($\Sigma ET_m\approx368$ mm). Find the relative and absolute yield, the loss, and compare the value of the water saved (0.12 per m$^3$) with the grain lost (250 per Mg).}\sol $Y/Y_m=1-K_y(1-ET_a/ET_m)=1-1.10\times0.20=0.780$ \ans $Y=0.780\times15.0=11.70$ Mg\,ha$^{-1}$ \ans Loss $=3.30$ Mg\,ha$^{-1}$, worth 825 per ha \ans Water saved $=0.2\times3680=736$ m$^3$\,ha$^{-1}$, worth 88 per ha \ans The grain lost is worth more than the water saved: the deficit does not pay at these prices.

}{}
\IfFileExists{extra/ex05.tex}{\backhead{EXERCISES}{khred}
\begin{exlist}
\item Find $e_s$ for $T_{max}=28$, $T_{min}=14$ and $e_a$ for RH = 60 \%.\exans{2.689 kPa; 1.614 kPa}
\item Compute $ET_0$ for $T_m=20$, $T_x=26$, $T_n=14$, $R_n=11$, $u_2=1.5$, RH = 65 \%, altitude 100 m.\exans{3.73 mm d$^{-1}$}
\item Find $K_c$ for wheat on days 10, 35 and 100.\exans{0.40, 0.77, 1.00}
\item Find $ET_c$ of tomato at day 50 for $ET_0=6.0$.\exans{5.10 mm d$^{-1}$}
\item If $U=7$ mm and $E_p=3.5$ mm\,d$^{-1}$, how long does stage 1 last?\exans{2 d}
\item Cumulative stage-2 evaporation after 6 days with $\beta=3.0$ mm\,d$^{-1/2}$.\exans{7.3 mm}
\item $ET_c=480$, $ET_a=400$, $K_y=1.1$, $Y_m=9$ Mg\,ha$^{-1}$: find $Y_a$.\exans{7.35 Mg ha$^{-1}$}
\item A field of 12 ha has NIR = 340 mm and uses drip ($E_a=0.9$). Give the volume.\exans{45333 m$^3$}
\item Find the stress coefficient $K_s$ when TAW = 120 mm, $p=0.5$, $D_r=84$ mm.\exans{0.60}
\item Explain why $K_{c,end}$ is low for grain crops and high for fresh vegetables.\exans{grain dries to maturity; fresh crops stay green}
\end{exlist}
}{}
\IfFileExists{extra/dyk05.tex}{\dykband{DO YOU KNOW ?}
\begin{exlist}
\item What is the difference between $ET_0$, $ET_c$ and $ET_a$?
\item Name the two terms of the Penman--Monteith equation and when each dominates.
\item Describe the four stages of the $K_c$ curve.
\item What limits evaporation in stage 1 and in stage 2?
\item State the Doorenbos--Kassam yield relation and the meaning of $K_y$.
\item Why is the effective rainfall less than the rainfall?
\item How does a mulch change stage 1?
\item Why is a drip system more water-saving than a sprinkler system on the same crop?
\end{exlist}
}{}
\IfFileExists{extra/obj05.tex}{\dykband{OBJECTIVE TYPE QUESTIONS}
\begin{exlist}
\item $ET_c$ is\opt{$K_cET_0$}{$ET_0/K_c$}{$K_c+ET_0$}{$K_cK_s$}
\item The reference crop is a\opt{well-watered grass of 0.12 m}{bare soil}{maize}{forest}
\item Stage 1 of soil evaporation is limited by\opt{soil conductivity}{energy}{roots}{salt}
\item Stage 2 evaporation falls as\opt{the square root of time}{the square of time}{a constant}{a step}
\item $K_y>1$ means the crop is\opt{tolerant of stress}{sensitive to stress}{dead}{a tree only}
\item The development stage has $K_c$\opt{constant}{rising linearly}{falling}{zero}
\item The effective rainfall is\opt{greater than the rain}{part of the rain that the crop can use}{the runoff}{the drainage}
\item Sprinkler efficiency of 0.75 means the gross requirement is NIR\opt{$\times0.75$}{$/0.75$}{$-0.75$}{$+0.75$}
\item The psychrometric constant is about\opt{0.067 kPa/$^\circ$C}{6.7 kPa/$^\circ$C}{67 kPa/$^\circ$C}{0.0067 kPa/$^\circ$C}
\item In deficit irrigation the crop receives\opt{more than $ET_c$}{less than $ET_c$}{no water}{only rain}
\end{exlist}
\objkey{1.~(a)\quad 2.~(a)\quad 3.~(b)\quad 4.~(a)\quad 5.~(b)\quad 6.~(b)\quad 7.~(b)\quad 8.~(b)\quad 9.~(a)\quad 10.~(b)}
}{}
